\documentclass[journal]{IEEEtran}

% ---------- Packages ----------
\usepackage{amsmath,amssymb,amsfonts,amsthm,mathtools}
\usepackage{graphicx}
\graphicspath{{output/pdf/}}
\usepackage{booktabs}
\usepackage{cite}
\usepackage{url}
\usepackage{microtype}
\usepackage[caption=false,font=footnotesize]{subfig}
\usepackage{mathrsfs}
\usepackage{hyperref}
\hypersetup{hidelinks}
\usepackage{xcolor}
\colorlet{blue}{black}
\usepackage[shortlabels]{enumitem}
\setlist[enumerate]{nosep,leftmargin=*}

\allowdisplaybreaks
\renewcommand{\topfraction}{0.90}
\renewcommand{\dbltopfraction}{0.95}
\renewcommand{\textfraction}{0.08}
\renewcommand{\floatpagefraction}{0.80}
\renewcommand{\dblfloatpagefraction}{0.99}
\setcounter{topnumber}{2}
\setcounter{dbltopnumber}{2}
\setlength{\dbltextfloatsep}{4pt plus 2pt minus 2pt}
\setlength{\textfloatsep}{4pt plus 2pt minus 2pt}

% ---------- Theorem Environments ----------
\newtheorem{theorem}{Theorem}
\newtheorem{proposition}{Proposition}
\newtheorem{corollary}{Corollary}
\newtheorem{lemma}{Lemma}
\newtheorem{assumption}{Assumption}
\newtheorem{definition}{Definition}
\theoremstyle{remark}
\newtheorem{remark}{Remark}

% ---------- Macros ----------
\newcommand{\R}{\mathbb{R}}
\newcommand{\E}{\mathbb{E}}
\newcommand{\inner}[2]{\left\langle #1,#2\right\rangle}
\newcommand{\norm}[1]{\left\lVert #1\right\rVert}
\newcommand{\Id}{\mathsf{I}}
\newcommand{\diag}{\operatorname{diag}}
\newcommand{\sym}{\operatorname{sym}}
\newcommand{\skewop}{\operatorname{skew}}
\newcommand{\cD}{\mathcal{D}}
\newcommand{\cN}{\mathcal{N}}

\newcommand{\bz}{\boldsymbol{z}}
\newcommand{\bu}{\boldsymbol{u}}
\newcommand{\btheta}{\boldsymbol{\theta}}
\newcommand{\bpsi}{\boldsymbol{\psi}}
\newcommand{\bzero}{\boldsymbol{0}}
\newcommand{\bH}{\boldsymbol{H}}
\newcommand{\bW}{\boldsymbol{W}}
\newcommand{\bS}{\boldsymbol{S}}
\newcommand{\bA}{\boldsymbol{A}}
\newcommand{\bM}{\boldsymbol{M}}

\newcommand{\bF}{\boldsymbol{F}}
\newcommand{\bG}{\boldsymbol{G}}

\newcommand{\hatF}{\widehat{\boldsymbol{F}}}
\newcommand{\hatG}{\widehat{\boldsymbol{G}}}

\newcommand{\biasF}{\boldsymbol{b}^{F}}
\newcommand{\biasG}{\boldsymbol{b}^{G}}
\newcommand{\noiseF}{\boldsymbol{\xi}^{F}}
\newcommand{\noiseG}{\boldsymbol{\xi}^{G}}

\newcommand{\Vk}{\mathcal{V}}
\newcommand{\cE}{\mathcal{E}}
\newcommand{\cP}{\mathcal{P}}

% ---------- Title ----------
\title{Saddle-Field Policy Optimization for Zero-Sum Markov Games: Adaptive Lyapunov-Drift Control with Finite-Time Stationarity Guarantees}

\author{Jialin~Zhuang,~\IEEEmembership{Member,~IEEE,}
	and~Vincent~K.~N.~Lau,~\IEEEmembership{Fellow,~IEEE}%
	\thanks{J.~Zhuang and V.~Lau are with the Department of Electronic and Computer Engineering, The Hong Kong University of Science and Technology, Clear Water Bay, Hong Kong Special Administrative Region, China (e-mail: jzhuangag@connect.ust.hk; eeknlau@ece.ust.hk).
		This work was supported by the Research Grants Council (RGC), Hong Kong, under the Areas of Excellence Scheme Grant AoE/E-101/23-N and AoE/E-601/22-R.
		(\em{Corresponding author: Vincent K. N. LAU}.)}
}

\begin{document}
	\maketitle
	\setlength{\abovedisplayskip}{4pt plus 2pt minus 2pt}
	\setlength{\belowdisplayskip}{4pt plus 2pt minus 2pt}
	\setlength{\abovedisplayshortskip}{1pt plus 1pt}
	\setlength{\belowdisplayshortskip}{4pt plus 2pt minus 2pt}
	
	% ============================================================
\begin{abstract}
Two-player zero-sum Markov games provide a framework for robust reinforcement learning and competitive policy optimization, but rotational coupling in their saddle fields can hinder first-order updates.
We combine the saddle-field direction $\bF(\bz)$ with the Jacobian-induced curvature direction $\bG(\bz)=\nabla\bF(\bz)\bF(\bz)$ and use a symmetric--antisymmetric Jacobian decomposition to characterize exactly when curvature decreases field energy.
A geometry-aware controller assigns separate step sizes to the two directions by minimizing a second-order model of the one-step change in a composite Lyapunov function, yielding a closed-form box-constrained quadratic program (QP).
For biased stochastic policy oracles, local geometric coverage gives an $O(1/K)$ bound on the time-average expected stationarity residual over $K$ joint updates, plus an explicit error floor.
Positive-semidefinite symmetric action and nonsingularity of the metric-transformed Jacobian imply neighborhood coverage, including nonsingular pure rotation; under a local policy-space error bound, the result further controls the unregularized Nash gap and worst-case-return loss.
Across four neural population-oracle games, the controller raises final mean worst-case return and lowers exploitability relative to its field-only ablation.
Finite-trajectory experiments compare six methods while sweeping attack strength in four neural-policy games; the controller attains the lowest mean exploitability at all 28 environment--severity settings and the strongest overall mean unregularized worst-case-return profile.
\end{abstract}
	
	\begin{IEEEkeywords}
		Zero-sum Markov games, robust reinforcement learning, stochastic saddle fields, Lyapunov drift, adaptive step size.
	\end{IEEEkeywords}
	
	% ============================================================
	\section{Introduction}\label{sec:intro}
	% ============================================================
	\IEEEPARstart{A}{two-player} zero-sum Markov game describes a sequential interaction in which two decision makers optimize opposing objectives \cite{Shapley1953StochasticGames,Littman1994MarkovGames}.
	The state records the information needed for the next decision, such as a controlled system's physical configuration, operating condition, or interaction history.
	Each player follows a policy that maps the current state to a distribution over its available actions.
	Their joint action determines both the instantaneous payoff and the next-state distribution, and the players respectively maximize and minimize the same expected cumulative return.
	Key applications include robust control under uncertain dynamics, adversarial reinforcement learning (RL), and competitive multi-agent decision making, where the maximizing policy selects control or decision actions and the minimizing policy represents uncertainty, disturbances, or an adaptive opponent \cite{SuttonBarto2018,MorimotoDoya2005RobustRL,Iyengar2005RobustDP,NilimElGhaoui2005,Pinto2017RARL,Kamalaruban2020LangevinRARL}.
	
	Policy optimization in a zero-sum Markov game is more involved than in a single-agent Markov decision process (MDP).
	An MDP optimizes one policy under fixed reward and transition mechanisms, whereas both policies in a Markov game adapt and jointly determine the state distribution, return, and policy gradients.
	An update by either player therefore changes the objective encountered by the other and produces coupled ascent--descent dynamics rather than a single policy-improvement direction.
	Neural policy parameterization further yields a generally nonconvex--nonconcave objective, while trajectory sampling introduces noise and bias into its gradient observations.
	The Jacobian of the joint policy dynamics can also contain a large antisymmetric component that rotates the update around a stationary point.
	Consequently, simultaneous gradient descent--ascent (GDA) or single timescale primal-dual updates can cycle or diverge even with small step sizes; related stability effects arise in linear--quadratic robust RL \cite{ZhangHuBasar2020StabilityRARL}.
	
	Classical dynamic-programming and value-based methods solve finite games through Bellman--Shapley operators or minimax value updates \cite{Shapley1953StochasticGames,Littman1994MarkovGames,Iyengar2005RobustDP,NilimElGhaoui2005}.
	However, their tabular implementations suffer from the curse of dimensionality because the computational and storage requirements grow with the state and joint-action spaces.
		This limitation motivates parameterized, trajectory-based policy optimization for larger games.
	Policy-gradient, actor--critic, optimistic, and natural-gradient schemes update the two policies directly and provide convergence or sample-complexity guarantees under different information and regularity models \cite{Sutton2000PolicyGradient,Wei2021LastIterate,Zhao2022EfficientPO,Zhang2022UnifiedPO,Alacaoglu2022NACMG}.
	These methods establish policy optimization as a scalable approach to Markov games, while the interaction between simultaneous player updates remains central to their stability.
	Competitive policy optimization incorporates this interaction through a local bilinear model of the two objectives \cite{Prajapat2021CoPO}.
	
	Proximal and extrapolation methods provide a complementary route for controlling coupled saddle dynamics.
	The proximal point method (PPM), the extragradient method (EGM), Mirror-Prox, and policy extragradient introduce corrections that can suppress rotational motion under suitable regularity and step-size conditions \cite{Rockafellar1976PPM,Korpelevich1976Extragradient,Nemirovski2004MirrorProx,Mokhtari2020UnifiedEG,CenWeiChi2024FastPolicyEG}.
	Local expansions identify a Jacobian-induced term as the leading correction shared by PPM and EGM beyond the GDA field direction \cite{Lu2022ODE,Mokhtari2020UnifiedEG}.
	Variable-step-size extragradient methods tune the extrapolation and update step sizes from observed operator information \cite{Hsieh2020DSEG,Antonakopoulos2021AdaProx}, while differentiable-game methods construct Hamiltonian, consensus, or symplectic adjustments from the game Jacobian \cite{Balduzzi2018Mechanics,Letcher2019DifferentiableGameMechanics,Mescheder2018GANConvergence,Loizou2020SHGD}.
	These developments expose curvature as a mechanism for controlling rotation, and they leave a concrete policy-optimization question: how should the field and curvature contributions be balanced as the local game geometry and stochastic oracle quality change along training?
	
	In this paper, we develop a geometry-adaptive stochastic policy optimizer for zero-sum Markov games that answers this question through two search directions with independently adaptive step sizes.
	The step sizes play a critical role in determining convergence and its rate in stochastic games.
	The saddle-field analysis supplies a first-order field direction and a Jacobian-induced curvature direction, and an exact symmetric--antisymmetric decomposition characterizes when the curvature direction decreases field energy.
	At each joint policy update, stochastic observations of these directions are combined with two independently selected nonnegative step sizes.
	A composite Lyapunov function measures both parameter-space stationarity and the policy-space improvement available through unilateral deviations.
	The controller selects the two step sizes by minimizing a second-order local model of the conditional Lyapunov drift over a two-dimensional box.
	Thus, the saddle geometry determines the search directions, whereas the Lyapunov-drift model determines their adaptive step sizes.
	The construction extends Lyapunov-drift step-size design for stochastic optimization \cite{Zhang2025AdaptivePolyak,Yuan2025StepSizeSDE} to a bi-direction saddle-field update and yields a closed-form box-constrained quadratic program (QP).
	The main contributions of this paper are summarized as follows.
	\begin{itemize}[leftmargin=*,nosep]
\item \textbf{Adaptive Bi-Directional Policy Optimization.}
		We propose a generic single timescale policy optimization algorithm with dynamic and state-dependent step sizes to balance the field and curvature directions.
		The step sizes are dynamically determined by minimizing a composite Lyapunov drift function, which combines first-order stationarity with regularized policy-space performance on a common scale. The drift minimization problem is equivalent to a two-variable box QP with a closed-form step-size solution.
		Specifically, the QP dominates every feasible field-only, curvature-only, or prescribed-coupling update in the same box, with a reduced-gradient condition for strict improvement.
\item \textbf{Geometric Characterization of the Curvature Direction.}
		We show that the curvature correction shared by PPM and EGM decreases field energy exactly when the antisymmetric Jacobian action exceeds the symmetric action.
		The resulting criterion identifies the local rotational geometry that motivates adaptive curvature use.
\item \textbf{Finite-Time and Policy-Space Guarantees.}
		For biased stochastic policy oracles, bi-direction geometric coverage yields an $O(1/K)$ transient in the time-average expected stationarity residual over $K$ joint policy updates, together with an explicit stochastic error floor.
		We establish neighborhood coverage from a positive-semidefinite symmetric Jacobian part and nonsingularity of the metric-transformed Jacobian at the target, including the nonsingular pure-rotation case.
		Under a local policy-space error bound, the stationarity guarantee further yields bounds on the ergodic Nash gap and the maximizing player's worst-case-return loss.
\item \textbf{Evidence Across Population and Finite-Trajectory Regimes.}
Neural population-oracle studies isolate the benefit of curvature under exact expected directions and nonlinear policy parameterization.
Finite-trajectory attack-strength sweeps across four neural-policy games compare six methods and yield the lowest mean exploitability at all 28 settings together with the strongest overall mean worst-case-return profile.
	\end{itemize}
	
	The remainder of the paper develops the argument as follows.
	Section~\ref{sec:prelim} formulates the parameterized Markov game, derives the two search directions, and introduces their stochastic recursion.
	Section~\ref{sec:drift} constructs the Lyapunov-drift step-size controller, Section~\ref{sec:conv} establishes the finite-time and policy-space guarantees, Section~\ref{sec:impl} gives the practical implementation, and Section~\ref{sec:experiments} presents the experiments.
	
	% ============================================================
	\section{Zero-Sum Markov Games and Stochastic Bi-Direction Policy Dynamics}\label{sec:prelim}
	
	% ============================================================
	
	\subsection{Zero-Sum Markov Game and Saddle-Field Formulation}\label{subsec:formulation}
	
	\textcolor{blue}{We consider an infinite-horizon discounted two-player zero-sum Markov game with state space $\mathcal{S}$, finite maximizing- and minimizing-player action spaces $\mathcal{A}$ and $\mathcal{B}$, and discount factor $\alpha\in(0,1)$.}
	\textcolor{blue}{For $s\in\mathcal{S}$, $a\in\mathcal{A}$, and $b\in\mathcal{B}$, the state-transition kernel is $P(\cdot\mid s,a,b)\triangleq\Pr(s_{t+1}\in\cdot\mid s_t=s,a_t=a,b_t=b)$, the bounded stage-reward function for the maximizing player is $r:\mathcal{S}\times\mathcal{A}\times\mathcal{B}\to\mathbb{R}$, and the initial-state distribution is $\rho(\cdot)\triangleq\Pr(s_0\in\cdot)$.}
	For a maximizing-player policy $\pi$ and a minimizing-player policy $\nu$, the maximizing player's discounted return is
	\[
	J(\pi,\nu)
	\triangleq
	\E_{\rho,\pi,\nu}
	\!\left[
	\textstyle\sum_{t=0}^{\infty}\alpha^{t}r(s_{t},a_{t},b_{t})
	\right].
	\]
	\textcolor{blue}{Here $s_t\in\mathcal{S}$, $a_t\in\mathcal{A}$, and $b_t\in\mathcal{B}$ denote the state, maximizing-player action, and minimizing-player action at time $t$, respectively.}
	\textcolor{blue}{The expectation $\mathbb{E}_{\rho,\pi,\nu}$ is over trajectories generated by $s_0\sim\rho$, $a_t\sim\pi(\cdot\mid s_t)$, $b_t\sim\nu(\cdot\mid s_t)$, and $s_{t+1}\sim P(\cdot\mid s_t,a_t,b_t)$.}
	% The expectation $\E_{\rho,\pi,\nu}$ is taken over the trajectory induced by the initial-state distribution, the two policies, and the transition kernel.
	The policy-space saddle objective is $\sup_{\pi}\inf_{\nu} J(\pi,\nu)$.
	Define the discounted causal entropies
	\begin{align*}
		\mathsf H_{\alpha}^{\pi}(\pi,\nu)
		&\triangleq \E_{\rho,\pi,\nu}\!\left[\sum_{t\ge0}\alpha^t
		\bigl(-\log\pi(a_t\mid s_t)\bigr)\right],\\
		\mathsf H_{\alpha}^{\nu}(\pi,\nu)
		&\triangleq \E_{\rho,\pi,\nu}\!\left[\sum_{t\ge0}\alpha^t
		\bigl(-\log\nu(b_t\mid s_t)\bigr)\right],
	\end{align*}
	with $0\log0=0$.
	For entropy weight $\varepsilon\ge0$, the corresponding regularized return is
	\begin{equation}\label{eq:entropic_return}
		J_{\varepsilon}(\pi,\nu)
		\triangleq J(\pi,\nu)
		+\varepsilon\mathsf H_{\alpha}^{\pi}(\pi,\nu)
		-\varepsilon\mathsf H_{\alpha}^{\nu}(\pi,\nu).
	\end{equation}
	\textcolor{blue}{The opposite signs preserve the max--min convention, while the entropy terms encourage stochastic policies for both players and smooth policy-space best-response values when the active action changes \cite{CenWeiChi2024FastPolicyEG}.}
	Setting $\varepsilon=0$ recovers the unregularized return $J$.
	We use $J_\varepsilon$ to construct the policy-update field and retain $J$ for the final unregularized performance measures.
	\textcolor{blue}{This smoothing supplies the differentiability required by the local step-size model in Section~\ref{sec:drift}.}
	\textcolor{blue}{The bridge from the regularized guarantee to the unregularized Nash gap and worst-case return is established in Section~\ref{subsec:performance_guarantee}.}
	In large finite or continuous state spaces, we adopt the standard parameterized formulation because exact policy-space optimization is generally intractable.
	The maximizing-player policy is represented by $\pi_{\btheta}(a\mid s)$ with $\btheta\in\R^{d_{\theta}}$, and the minimizing-player policy by $\nu_{\bpsi}(b\mid s)$ with $\bpsi\in\R^{d_{\psi}}$, giving
	\[
	\begin{aligned}
		J_{\varepsilon}(\btheta,\bpsi)
		&\triangleq
		J_{\varepsilon}(\pi_{\btheta},\nu_{\bpsi}),\\
		\max_{\btheta}\min_{\bpsi}\quad
		&J_{\varepsilon}(\btheta,\bpsi).
	\end{aligned}
	\]
	Write $\bz\triangleq(\btheta,\bpsi)\in\R^d$, where $d\triangleq d_\theta+d_\psi$.
	The joint policy ascent--descent dynamics define the saddle field
	\begin{equation}\label{eq:Fdef}
		\bF(\bz)
		\triangleq
		\begin{bmatrix}
			-\nabla_{\btheta}J_{\varepsilon}(\btheta,\bpsi)\\[0.2em]
			\phantom{-}\nabla_{\bpsi}J_{\varepsilon}(\btheta,\bpsi)
		\end{bmatrix}.
	\end{equation}
	The first-order stationary set is $\mathcal Z^{\star}\triangleq\{\bz:\bF(\bz)=\bzero\}$, where $\bzero$ is the zero vector.
	Every $\bz^\star\in\mathcal Z^\star$ satisfies the unconstrained first-order Karush--Kuhn--Tucker (KKT) conditions for the parameterized objective $J_\varepsilon$.
	Under neural policy parameterization, $\mathcal Z^\star$ is the parameter-space first-order stationary set.
	This distinction matters because a nonconvex or noninjective policy parameterization can satisfy $\bF(\bz)=\bzero$ even when a player can still improve by taking a best response in the full policy space.
	The policy-space results in Section~\ref{subsec:performance_guarantee} use the Nash gap to measure this remaining unilateral-improvement incentive and additionally bound the maximizing player's worst-case-return loss.
	
	\subsection{Classical Saddle Updates and the Curvature Direction}\label{subsec:saddle_field}
	
	The saddle-field formulation exposes the rotational component that can destabilize first-order updates.
	To understand how established saddle methods counter this effect, we examine the Jacobian-based correction generated by PPM and EGM.
	For a scalar step size $s>0$, let $\bz_{k+1}$ denote the next joint policy parameter.
	The representative GDA~\cite{Mokhtari2020UnifiedEG}, PPM~\cite{Rockafellar1976PPM}, and EGM~\cite{Korpelevich1976Extragradient} updates are
	\begin{equation}\label{eq:classical_updates}
		\begin{aligned}
			\text{\normalfont GDA~\cite{Mokhtari2020UnifiedEG}:}\quad&
			\bz_{k+1}^{\rm GDA}
			=\bz_k-s\bF(\bz_k),\\
			\text{\normalfont PPM~\cite{Rockafellar1976PPM}:}\quad&
			\bz_{k+1}^{\rm PPM}
			=\bz_k-s\bF(\bz_{k+1}^{\rm PPM}),\\
			\text{\normalfont EGM~\cite{Korpelevich1976Extragradient}:}\quad&
			\widetilde{\bz}_k
			=\bz_k-s\bF(\bz_k),\quad
			\bz_{k+1}^{\rm EGM}
			=\bz_k-s\bF(\widetilde{\bz}_k).
		\end{aligned}
	\end{equation}
	where $\widetilde{\bz}_k$ is the EGM extrapolation point.
	GDA is explicit, PPM is implicit, and EGM uses an extrapolated field query.
	
	\begin{figure*}[!t]
		\centering
		\includegraphics[width=\textwidth]{fig_motivation_geometry_abd_wide.pdf}
		\caption{Geometry of bi-direction saddle updates.
			(a) GDA spirals outward, whereas PPM and EGM contract at the displayed step size.
			(b) On the bilinear game, the first-order field direction is tangential to the residual contours, whereas the shared second-order correction points inward.
			(c) The curvature-to-field step-size ratio controls contraction and residual rotation.}
		\label{fig:motivation_geometry}
	\end{figure*}
	
	The following assumption states the local regularity required by the template expansion and one-step analysis.
	
	\begin{assumption}[Local smoothness]\label{ass:smooth}
		There exists an open set $\cD\subset\R^{d}$ such that:
		(a) all iterates and probing points used in the analysis lie in $\cD$;
		(b) $J_{\varepsilon}:\R^{d}\to\R$ is four times continuously differentiable ($C^{4}$) on $\cD$;
		(c) all derivatives of $J_{\varepsilon}$ up to order four are bounded on $\cD$.
	\end{assumption}
	
	Conditions~(b)--(c) imply that $\bF$ is three times continuously differentiable, i.e. $\bF\in C^{3}(\cD;\R^{d})$; after prescribing an implementable region $\mathcal T\subset\cD$, the safeguards of Section~\ref{sec:impl} can maintain Condition~(a).
	Taylor expansion of the field uses its Jacobian $\nabla\bF$ and second-derivative tensor $\nabla^2\bF$, both taken with respect to $\bz$.
	The remainder bounds use the Euclidean norm $\norm{\cdot}$ and its induced operator and tensor norms.
	
	\begin{proposition}[Discrete-time second-order template expansions]\label{prop:DTA}
		Suppose Assumption~\ref{ass:smooth} holds.
		Fix $k\ge0$ and $\bz_k\in\cD$.
		There exists $\bar s_k>0$ such that, for every $s\in(0,\bar s_k]$, the PPM equation has a unique local solution, all displayed update and extrapolation points lie in $\cD$, and
		\begin{align*}
			\textnormal{GDA:}\quad \bz_{k+1}^{\rm GDA}
			&=\bz_k-s\bF(\bz_k),\\
			\textnormal{PPM:}\quad \bz_{k+1}^{\rm PPM}
			&=\bz_k-s\bF(\bz_k)+s^{2}\nabla\bF(\bz_k)\bF(\bz_k)\\
			&\quad{}+R^{\mathrm{PPM}}_{s}(\bz_k),\\
			\textnormal{EGM:}\quad \bz_{k+1}^{\rm EGM}
			&=\bz_k-s\bF(\bz_k)+s^{2}\nabla\bF(\bz_k)\bF(\bz_k)\\
			&\quad{}+R^{\mathrm{EGM}}_{s}(\bz_k).
		\end{align*}
		The remainders satisfy $\norm{R^{\mathrm{PPM}}_{s}(\bz_k)}\le C_{\mathrm{PPM}}s^{3}$ and $\norm{R^{\mathrm{EGM}}_{s}(\bz_k)}\le C_{\mathrm{EGM}}s^{3}$ for finite constants depending only on local bounds of $\bF$, $\nabla\bF$, and $\nabla^{2}\bF$ around $\bz_k$.
	\end{proposition}
\textcolor{blue}{The proof, including local PPM solvability and the stated remainder bounds, is given in \hyperref[app:DTA]{Appendix~\ref*{app:proofs}, item~a)}.}
	
	Unlike GDA, PPM and EGM share the leading-order correction $+s^{2}\nabla\bF(\bz)\bF(\bz)$, which can suppress rotational saddle dynamics at suitable step sizes.
	This contrast motivates writing $\bA(\bz)\triangleq\nabla\bF(\bz)$ and isolating the curvature direction
	\begin{equation}\label{eq:curvature_direction}
		\bG(\bz)\triangleq\bA(\bz)\bF(\bz).
	\end{equation}
	The expansion identifies the curvature direction shared by PPM and EGM, and the next step is to characterize how this direction changes the local saddle dynamics.
	
	\subsection{Curvature Descent and Independent Step Sizes}\label{subsec:curvature_geometry}
	
	Proposition~\ref{prop:DTA} identifies $+\bG$ as the leading correction shared by PPM and EGM.
	Its effect on the stationarity residual is governed by the local Jacobian geometry, which motivates characterizing when $+\bG$ decreases the Euclidean field energy
	\[
	\phi(\bz)\triangleq\tfrac12\norm{\bF(\bz)}^2,
	\]
	which is nonnegative and vanishes exactly on $\mathcal Z^\star$.
	
	Fig.~\ref{fig:motivation_geometry} first illustrates the two directions on the bilinear max--min saddle $L(x,y)=-xy$.
	Its field is $\bF=J_0\bz$, where
	\[
	J_0
	\triangleq
	\begin{bmatrix}
		0&1\\
		-1&0
	\end{bmatrix}.
	\]
	Since $J_0^2=-\Id_2$, where $\Id_2$ is the $2\times2$ identity matrix, the field direction $-\bF$ is tangent to the level sets of $\phi$, whereas $+\bG=J_0^2\bz=-\bz=-\nabla\phi$ points inward.
	This geometry explains why GDA rotates while PPM and EGM contract at the displayed step size in panel (a), and panel (b) shows the two directional roles directly.
	Panel (c) further shows that changing the curvature-to-field step-size ratio changes both contraction and residual rotation, motivating separate selection of the two step sizes.
	
	To extend this observation beyond the bilinear game, fix $\bz\in\cD$, write $\bA=\nabla\bF(\bz)$, and decompose the Jacobian as
	\[
	\begin{aligned}
		\bA&=\bS+\bW,\\
		\bS&=\sym(\bA)\triangleq\tfrac12(\bA+\bA^\top),\\
		\bW&=\skewop(\bA)\triangleq\tfrac12(\bA-\bA^\top).
	\end{aligned}
	\]
	The symmetric part $\bS$ captures nonrotational dynamics, whereas the antisymmetric part $\bW$ captures rotational dynamics \cite{Balduzzi2018Mechanics}.
	Using $\inner{\cdot}{\cdot}$ for the standard Euclidean inner product, the exact directional derivatives give the following criterion.
	
	\begin{theorem}[Skew-dominance characterization of curvature descent]\label{thm:skew_dominance}
		Suppose Assumption~\ref{ass:smooth} holds.
		For every $\bz\in\cD$ with $\bF(\bz)\neq\bzero$, the field and curvature directions satisfy
		\begin{align}
			\inner{\nabla\phi(\bz)}{-\bF(\bz)}
			&=-\bF(\bz)^\top\bS\bF(\bz),
			\label{eq:dphi_F}\\
			\inner{\nabla\phi(\bz)}{\bG(\bz)}
			&=\norm{\bS\bF(\bz)}^2-\norm{\bW\bF(\bz)}^2.
			\label{eq:dphi_G}
		\end{align}
		Consequently, $+\bG(\bz)$ strictly decreases $\phi$ if and only if
		\begin{equation}\label{eq:skew_dominance}
			\norm{\bW\bF(\bz)}
			>
			\norm{\bS\bF(\bz)}.
		\end{equation}
	\end{theorem}
	
	Eq.~\eqref{eq:dphi_F} shows that the symmetric Jacobian action governs the first-order effect of the field direction.
	For a purely rotational field, $\bS=0$ makes $-\bF$ tangential, while $\bW\bF\neq\bzero$ makes $+\bG$ a strict descent direction.
	Conversely, when $\bW=0$, Eq.~\eqref{eq:dphi_G} shows that $+\bG$ preserves or increases the field energy to first order.
	For a two-dimensional normal block $\bA=\mu\Id_2+\sigma J_0$, the criterion reduces to $\lvert\sigma\rvert>\lvert\mu\rvert$, directly comparing rotational and radial action.
	
	Because this balance changes with the local Jacobian geometry, we assign independent nonnegative step sizes $\beta$ and $\gamma$ to the two population directions:
	\begin{equation}\label{eq:population_two_direction}
		\bz_{k+1}
		=
		\bz_k-\beta\bF(\bz_k)+\gamma\bG(\bz_k),
		\quad
		\beta,\gamma\ge0.
	\end{equation}
	The pair $(\beta,\gamma)=(s,0)$ recovers GDA, whereas $(\beta,\gamma)=(s,s^2)$ matches the PPM and EGM templates through second order.
	Eq.~\eqref{eq:population_two_direction} therefore retains the classical search directions while allowing their step sizes to adapt independently to the local game geometry.
	
	\subsection{Stochastic Bi-Direction Recursion and Oracle Model}\label{subsec:oracle_model}
	
	Eq.~\eqref{eq:population_two_direction} defines the joint-update family underlying the proposed recursion: the field direction retains the first-order GDA action, while the curvature direction generalizes the leading correction shared by PPM and EGM.
	In trajectory-based policy optimization, both population directions are observed through stochastic policy oracles.
	At joint update $k$, let $\hatF_k$ and $\hatG_k$ denote these observations and let $\beta_k,\gamma_k\ge0$ denote their step sizes.
	The resulting implementable bi-direction recursion is
	\begin{equation}\label{eq:update}
		\bz_{k+1}
		=
		\bz_{k}-\beta_{k}\hatF_{k}+\gamma_{k}\hatG_{k},
	\end{equation}
	Section~\ref{subsec:trajectory_oracles} constructs both observations from finite trajectory batches.
	For the stochastic analysis, let $\{\mathscr F_k\}_{k\ge0}$ be the filtration generated by the algorithmic history revealed before the queries at joint update $k$.
	
	\begin{assumption}[Biased stochastic oracles]\label{ass:oracles}
		For every $k$, let $\biasF_k,\biasG_k$ be $\mathscr F_k$-measurable oracle biases and let $\noiseF_k,\noiseG_k$ be conditionally zero-mean noise variables, i.e. $\E[\noiseF_k\mid\mathscr F_k]=\E[\noiseG_k\mid\mathscr F_k]=\bzero$.
		The observed directions decompose as
		\[
		\hatF_{k}
		=
		\bF(\bz_{k})+\biasF_{k}+\noiseF_{k},
		\quad
		\hatG_{k}
		=
		\bG(\bz_{k})+\biasG_{k}+\noiseG_{k},
		\]
		There exist finite $\mathscr{F}_{k}$-measurable scalars $\bar b^{F}_{k},\bar b^{G}_{k},\sigma_{2,k},\sigma_{3,k}\ge 0$ such that, almost surely,
		\begin{enumerate}[(a)]
			\item $\norm{\biasF_{k}}\le \bar b^{F}_{k}$ and $\norm{\biasG_{k}}\le \bar b^{G}_{k}$;
			\item $\E[\norm{\noiseF_{k}}^{2}+\norm{\noiseG_{k}}^{2} \mid\mathscr{F}_{k}]\le \sigma_{2,k}^{2}$;
			\item $\E[\norm{\hatF_{k}}^{3}+\norm{\hatG_{k}}^{3} \mid\mathscr{F}_{k}]\le \sigma_{3,k}^{3}$.
		\end{enumerate}
	\end{assumption}
	
	Here $\bar b_k^F,\bar b_k^G$ bound the conditional biases of the two population-direction estimates, while $\sigma_{2,k},\sigma_{3,k}$ bound their second noise moment and third oracle moment.
	The model permits mini-batch estimates as well as bounded errors from trajectory truncation, clipped surrogates, and approximate curvature evaluations.
	Eq.~\eqref{eq:update} specifies the search directions and their stochastic observations.
	The remaining design task is to determine their two step sizes.
	Independent selection of $\beta_k$ and $\gamma_k$ allows the two step sizes to respond separately to local rotation, policy-space performance, and oracle quality.
	The next section therefore constructs a composite Lyapunov function and selects $(\beta_k,\gamma_k)$ from its local one-step model.
	
	% ============================================================
	\section{Geometry-Aware Lyapunov-Drift Step-Size Selection}\label{sec:drift}
	% ============================================================
	Section~\ref{sec:prelim} determines the field and curvature directions of the stochastic policy update.
	The remaining design problem is to select their two step sizes from the current game geometry and oracle observations.
	Field energy measures parameter-space stationarity, while policy-space performance requires accounting for the improvement available through unilateral deviations.
	We place these quantities on a common Lyapunov scale, model its conditional one-step change, and select the two step sizes from an implementable local approximation.
	
	\subsection{Composite Lyapunov Function}\label{subsec:V_drift}
	The field energy measures parameter-space stationarity.
	The regularized Nash gap measures unilateral policy improvement and can distinguish parameter-stationary policies that retain positive exploitability.
	Because these quantities have different numerical scales, the following Lyapunov function combines them: field energy guides local optimization, while the regularized gap retains policy-space performance in the step-size decision.
	
	\textcolor{blue}{Let $\diag(\cdot)$ denote block-diagonal concatenation, set $\bW_{V}\triangleq\diag(\Id_{d_{\theta}},\eta_{\psi}\Id_{d_{\psi}})$ for a relative block weight $\eta_{\psi}>0$, and define $\cE_{F}(\bz)\triangleq\tfrac12\bF(\bz)^{\top}\bW_{V}\bF(\bz)$.}
	\textcolor{blue}{Its population decrease signals along the field and curvature directions are}
	\begin{equation}\label{eq:geometry_signals}
		\begin{aligned}
			\chi_F(\bz)
			&\triangleq-\inner{\nabla\cE_F(\bz)}{-\bF(\bz)}\\
			&=\tfrac12\bF^\top(\bW_V\bA+\bA^\top\bW_V)\bF,\\
			\chi_G(\bz)
			&\triangleq-\inner{\nabla\cE_F(\bz)}{\bG(\bz)}\\
			&=-\bF^\top\bW_V\bA^2\bF.
		\end{aligned}
	\end{equation}
	\textcolor{blue}{Positive values indicate first-order decrease of the weighted field energy along the corresponding direction.}
	\textcolor{blue}{For $\varepsilon>0$, write $J_\varepsilon=J_\varepsilon(\pi_{\btheta},\nu_{\bpsi})$ and define the policy-performance component}
	\[
	\begin{aligned}
		\cP(\bz)
		&\triangleq\operatorname{Gap}_{\varepsilon}(\pi_{\btheta},\nu_{\bpsi})\\
		&=\left[\sup_{\bar\pi}J_{\varepsilon}(\bar\pi,\nu_{\bpsi})-J_\varepsilon\right]
		+\left[J_\varepsilon-\inf_{\bar\nu}J_{\varepsilon}(\pi_{\btheta},\bar\nu)\right].
	\end{aligned}
	\]
	\textcolor{blue}{The extrema range over stationary randomized policies, so $\cP$ is the sum of the players' unilateral regularized improvements and vanishes when neither player can improve.}
	\textcolor{blue}{Fix normalizers $\kappa_{F},\kappa_{P}>0$ and weights $\lambda_{F}>0$, $\lambda_{P}\ge0$.}
	\begin{definition}[Composite Lyapunov function]\label{def:Lyap}
		\textcolor{blue}{The composite Lyapunov function is}
		\begin{equation}\label{eq:Vdef}
			\Vk(\bz)
			\triangleq
			\lambda_{F}\frac{\cE_{F}(\bz)}{\kappa_{F}}
			+
			\lambda_{P}\frac{\cP(\bz)}{\kappa_{P}}.
		\end{equation}
	\end{definition}
	\textcolor{blue}{Here $\bW_V$ is the metric weight, distinct from $\bW=\skewop(\nabla\bF)$ in Section~\ref{subsec:curvature_geometry}; $\cP\equiv0,\lambda_P=0$ give the field-energy-only controller.}
	
	For the recursion in Eq.~\eqref{eq:update}, define the realized one-step Lyapunov increment by $\mathscr I_k\triangleq\Vk(\bz_{k+1})-\Vk(\bz_k)$.
	Its conditional expectation
	\[
	\mathscr D_k\triangleq\E[\mathscr I_k\mid\mathscr F_k]
	\]
	is the Lyapunov drift at update $k$.
	The increment $\mathscr I_k$ is pathwise after the current queries are realized, whereas $\mathscr D_k$ averages it conditional on the pre-query history $\mathscr F_k$.
	Hence, $\mathscr D_k<0$ represents expected one-step reduction of the combined stationarity and policy-performance measure, which makes its local minimization the step-size design objective.
	
	Theorem~\ref{thm:skew_dominance} analyzes $+\bG$ for the field energy $\phi$, whereas Definition~\ref{def:Lyap} extends the same geometry to $\Vk$ so that the local drift model can balance curvature use against policy performance.
	
	Since $\cP\ge0$, $\bW_V$ is positive definite, and $\lambda_F>0$, Definition~\ref{def:Lyap} gives $\Vk\ge0$ and $\Vk(\bz)=0$ only if $\bz\in\mathcal Z^\star$.
	Under Assumption~\ref{ass:smooth} and $\cP\in C^3(\cD)$, the function $\Vk$ is $C^3$ with gradient $\nabla\Vk=(\lambda_F/\kappa_F)\nabla\cE_F +(\lambda_P/\kappa_P)\nabla\cP$, where $\nabla\cE_F=\nabla\bF^{\top}\bW_V\bF$.
	The weighted stationarity residual $\Phi(\bz)\triangleq\bF(\bz)^{\top}\bW_{V}\bF(\bz)$ satisfies $\Phi(\bz)\le2\kappa_{F}\Vk(\bz)/\lambda_{F}$.
	
	The composite Lyapunov decrease signals along the two population directions are
	\begin{equation}\label{eq:composite_signals}
		\begin{aligned}
			D_F(\bz)
			&\triangleq\inner{\nabla\Vk(\bz)}{\bF(\bz)}
			\\[-2pt]
			&=\frac{\lambda_F}{\kappa_F}\chi_F(\bz)
			+\frac{\lambda_P}{\kappa_P}\inner{\nabla\cP(\bz)}{\bF(\bz)},\\
			D_G(\bz)
			&\triangleq-\inner{\nabla\Vk(\bz)}{\bG(\bz)}
			\\[-2pt]
			&=\frac{\lambda_F}{\kappa_F}\chi_G(\bz)
			-\frac{\lambda_P}{\kappa_P}\inner{\nabla\cP(\bz)}{\bG(\bz)}.
		\end{aligned}
	\end{equation}
	Thus, $D_F>0$ and $D_G>0$ mean first-order decrease of $\Vk$ along $-\bF$ and $+\bG$, respectively.
	
	\subsection{Second-Order Local Drift Model}\label{subsec:two_direction}
	For a realized oracle pair and candidate step sizes $(\beta,\gamma)$, define the trial displacement $\Delta_{k}(\beta,\gamma)\triangleq -\beta\hatF_{k}+\gamma\hatG_{k}$ and the pathwise trial increment $\mathscr I_k(\beta,\gamma)\triangleq\Vk(\bz_k+\Delta_k(\beta,\gamma))-\Vk(\bz_k)$.
	The actual quantities satisfy $\mathscr I_k=\mathscr I_k(\beta_k,\gamma_k)$ and $\mathscr D_k=\E[\mathscr I_k(\beta_k,\gamma_k)\mid\mathscr F_k]$.
	Exact query-dependent minimization of $\mathscr D_k$ requires the full conditional oracle distribution and repeated evaluation of the nonlinear post-update function $\Vk(\bz_k+\Delta_k)$.
	A single query instead provides one realized oracle pair, so we minimize a second-order pathwise surrogate of $\mathscr I_k(\beta,\gamma)$ after observing $(\hatF_k,\hatG_k)$.
	The surrogate exposes the directional curvature along each search direction and their cross-curvature, which determine the step-size pair $(\beta_k,\gamma_k)$; Section~\ref{sec:conv} bounds its mismatch from the realized increment and converts the pathwise choice into an expected finite-time guarantee.
	
	\begin{assumption}[Local regularity along the iterates]\label{ass:local_reg}
		For every $k$ and every feasible $(\beta,\gamma)$ produced by the step-size routine, the segment $\{\bz_{k}+t\Delta_{k}(\beta,\gamma):t\in[0,1]\}\subset\cD$.
		There exist $\mathscr{F}_{k}$-measurable $G_{V,k},H_{V,k},M_{3,k}\ge 0$ that, simultaneously for every feasible pair and almost every realization of the current stochastic queries, satisfy $\norm{\nabla\Vk(\bz_{k})}\le G_{V,k}$, $\norm{\nabla^{2}\Vk(\bz_{k})}\le H_{V,k}$, and $\sup_{t\in[0,1]}\norm{\nabla^{3}\Vk(\bz_{k}+t\Delta_{k}(\beta,\gamma))} \le M_{3,k}$.
	\end{assumption}
	
	Under Assumption~\ref{ass:local_reg}, Taylor's theorem gives the following expansion for every feasible $(\beta,\gamma)$:
	\begin{equation}\label{eq:drift_taylor}
		\begin{aligned}
			&\Vk(\bz_{k}+\Delta_{k}(\beta,\gamma))-\Vk(\bz_{k}) \\
			={}&-\beta\inner{\nabla\Vk(\bz_k)}{\hatF_k}
			+\gamma\inner{\nabla\Vk(\bz_k)}{\hatG_k}\\
			&+\tfrac12\beta^2\hatF_k^\top\nabla^2\Vk(\bz_k)\hatF_k
			-\beta\gamma\hatF_k^\top\nabla^2\Vk(\bz_k)\hatG_k\\
			&+\tfrac12\gamma^2\hatG_k^\top\nabla^2\Vk(\bz_k)\hatG_k
			+\mathrm{rem}_{k}(\beta,\gamma),
		\end{aligned}
	\end{equation}
	The remainder satisfies $\lvert \mathrm{rem}_{k}(\beta,\gamma)\rvert\le\frac{M_{3,k}}{6}(\beta\norm{\hatF_{k}}+\gamma\norm{\hatG_{k}})^{3}$.
	
	Eq.~\eqref{eq:drift_taylor} is the pathwise second-order Taylor expansion for each realized oracle pair.
	Its linear and quadratic terms define the five scalar coefficients
	\begin{equation}\label{eq:drift_terms_def}
		\begin{aligned}
			e_{k}&\triangleq\inner{\nabla\Vk(\bz_{k})}{\hatF_{k}},
			\quad
			d_{k}\triangleq -\inner{\nabla\Vk(\bz_{k})}{\hatG_{k}},\\
			c_{k}&\triangleq\hatF_{k}^{\top}\nabla^{2}\Vk(\bz_{k})\hatF_{k},
			\quad
			a_{k}\triangleq\hatG_{k}^{\top}\nabla^{2}\Vk(\bz_{k})\hatG_{k},\\
			b_{k}&\triangleq\hatF_{k}^{\top}\nabla^{2}\Vk(\bz_{k})\hatG_{k}.
		\end{aligned}
	\end{equation}
	Eq.~\eqref{eq:drift_terms_def} is pathwise: it uses the realized oracle pair $(\hatF_k,\hatG_k)$ rather than the population directions $(\bF(\bz_k),\bG(\bz_k))$.
	Collecting these five coefficients gives the pathwise reference quadratic
	\begin{equation}\label{eq:qref}
		q_k(\beta,\gamma)
		\triangleq
		-\beta e_k-\gamma d_k
		+\tfrac12c_k\beta^2-b_k\beta\gamma
		+\tfrac12a_k\gamma^2.
	\end{equation}
	Thus, Eq.~\eqref{eq:drift_taylor} is $\mathscr I_k(\beta,\gamma)=q_k(\beta,\gamma)+\mathrm{rem}_k(\beta,\gamma)$.
	Here $e_k,d_k$ are first-order predicted decreases along $-\hatF_k,+\hatG_k$, $c_k,a_k$ are their directional curvatures, and $b_k$ is the cross-curvature.
	Assumptions~\ref{ass:local_reg} and~\ref{ass:oracles} respectively control the third-order remainder and the conditional oracle errors used in the expected-drift analysis.
	
	\subsection{Closed-Form Box-QP Step-Size Selection}\label{subsec:closed_form}
	Discarding the bounded third-order remainder in Eq.~\eqref{eq:drift_taylor} yields a bivariate quadratic model for selecting $(\beta,\gamma)$.
	Eq.~\eqref{eq:drift_terms_def} defines reference coefficients from $\nabla\Vk(\bz_k)$ and directional Hessian actions.
	We denote their evaluated counterparts returned by the local-model routine by $(\tilde e_k,\tilde d_k,\tilde c_k,\tilde a_k,\tilde b_k)$.
	Exact automatic differentiation gives $(\tilde e_k,\tilde d_k,\tilde c_k,\tilde a_k,\tilde b_k)=(e_k,d_k,c_k,a_k,b_k)$, whereas finite-difference or directional-secant implementations, possibly combined with moving averages, generally give estimates such as $\tilde e_k\approx e_k$ and $\tilde c_k\approx c_k$.
	The quadratic coefficient matrix assembled from the evaluated values is $\widetilde{\bH}_k\triangleq\left[\begin{smallmatrix}\tilde c_k&-\tilde b_k\\-\tilde b_k&\tilde a_k\end{smallmatrix}\right]$.
	Because local nonconvexity or evaluation error may make $\widetilde{\bH}_k$ indefinite or nearly singular, we stabilize it by a common spectral shift.
	Let $\lambda_{\min}(\cdot)$ denote the smallest eigenvalue and fix a positive-definiteness floor $\lambda_{\rm pd}>0$.
	Using $[u]_+\triangleq\max\{u,0\}$ for the positive part of a scalar $u$, the shift is
	\begin{equation}\label{eq:spectral_shift}
		\lambda_k=[\lambda_{\rm pd}-\lambda_{\min}(\widetilde{\bH}_k)]_+,
		\quad
		\widehat{\bH}_k=\widetilde{\bH}_k+\lambda_k\Id_2.
	\end{equation}
	Here $\lambda_k$ is the smallest nonnegative common diagonal correction that raises the minimum eigenvalue to at least $\lambda_{\rm pd}$.
	We denote the resulting numerically stabilized coefficients by hats: $\hat e_k=\tilde e_k$, $\hat d_k=\tilde d_k$, $\hat b_k=\tilde b_k$, and $(\hat c_k,\hat a_k)=(\tilde c_k+\lambda_k,\tilde a_k+\lambda_k)$.
	Equivalently, $\rho_{c,k}=\rho_{a,k}=\lambda_k\ge0$ add $\tfrac12\rho_{c,k}\beta^2+\tfrac12\rho_{a,k}\gamma^2$ to the tilde-coefficient model while leaving its linear and cross terms unchanged.
	The implemented model is therefore
	\begin{equation}\label{eq:qhat}
		\begin{aligned}
			\hat q_{k}(\beta,\gamma)
			&\triangleq
			-\beta\hat e_{k}-\gamma\hat d_{k}
			+\tfrac{1}{2}\hat c_{k}\beta^{2}-\hat b_{k}\beta\gamma\\
			&\quad+\tfrac{1}{2}\hat a_{k}\gamma^{2},
			\quad \textcolor{blue}{(\beta,\gamma)\in\R^{2}.}
		\end{aligned}
	\end{equation}
	Its stabilized Hessian is $\widehat{\bH}_k$, and $\hat\delta_k\triangleq\det\widehat{\bH}_k=\hat c_k\hat a_k-\hat b_k^2$.
	
	Eq.~\eqref{eq:spectral_shift} gives $\lambda_{\min}(\widehat{\bH}_k)=\max\{\lambda_{\min}(\widetilde{\bH}_k),\lambda_{\rm pd}\}\ge\lambda_{\rm pd}$.
	Thus, $\widehat{\bH}_k$ is positive definite, and its leading principal minors satisfy $\hat c_k>0$ and $\hat\delta_k>0$.
	\textcolor{blue}{For caps $\beta_{\max},\gamma_{\max}>0$, the controller solves}
	\begin{equation*}\tag{P1}\label{prob:box_qp}
		\textcolor{blue}{
		\begin{aligned}
		(\beta_k,\gamma_k)\in
		\arg\min_{\beta,\gamma}\quad &\hat q_k(\beta,\gamma)\\
		\mathrm{s.t.}\quad &0\le\beta\le\beta_{\max},\quad
		0\le\gamma\le\gamma_{\max}.
		\end{aligned}}
	\end{equation*}
	\textcolor{blue}{The decision variables $\beta$ and $\gamma$ are the field and curvature step sizes. The objective $\hat q_k$ is the stabilized second-order model of the one-step Lyapunov increment; minimizing it therefore selects the feasible update with the largest predicted local Lyapunov decrease. The nonnegativity constraints preserve the prescribed orientations $-\hatF_k$ and $+\hatG_k$, while the upper bounds limit the trial displacement to the local-model scale and provide finite caps for the safeguard in Section~\ref{subsec:safeguard}. To solve Problem~\eqref{prob:box_qp}, first consider the uncapped critical point}
	\begin{equation}\label{eq:unconstrained_pair}
		(\beta_{k}^{\rm uc},\gamma_{k}^{\rm uc})
		\triangleq
		\Bigl(\tfrac{\hat a_{k}\hat e_{k}+\hat b_{k}\hat d_{k}}{\hat\delta_{k}},\,
		\tfrac{\hat b_{k}\hat e_{k}+\hat c_{k}\hat d_{k}}{\hat\delta_{k}}\Bigr),
	\end{equation}
	and the boundary candidates are $\bigl([\hat e_{k}/\hat c_{k}]_{+},0\bigr)$, $\bigl(0,[\hat d_{k}/\hat a_{k}]_{+}\bigr)$, $(0,0)$.
	
	\begin{proposition}[Closed-form bi-direction surrogate controller]\label{prop:closed_form}
		\textcolor{blue}{Problem~\eqref{prob:box_qp} has a unique minimizer.}
		\textcolor{blue}{For its uncapped counterpart on $[0,\infty)^2$, the minimizer is given by Eq.~\eqref{eq:unconstrained_pair} when both components are nonnegative and otherwise by the boundary candidate with smallest $\hat q_k$.}
		When the unconstrained critical point belongs to $[0,\infty)^2$, the predicted decrease $\widehat\Delta_{k}^{\star}\triangleq -\hat q_{k}(\beta_{k}^{\star},\gamma_{k}^{\star})$ is
		\begin{equation}\label{eq:Delta_hat_interior}
			\widehat\Delta_{k}^{\star}=\frac{\hat a_{k}\hat e_{k}^{2}+2\hat b_{k}\hat e_{k}\hat d_{k}+\hat c_{k}\hat d_{k}^{2}}{2\hat\delta_{k}}.
		\end{equation}
		\textcolor{blue}{For the finite caps in Problem~\eqref{prob:box_qp}, the unique minimizer is the member with smallest objective value among the feasible unconstrained point and the four edge minimizers}
		\begin{align*}
			&\left(\left[\hat e_k/\hat c_k\right]_{[0,\beta_{\max}]},0\right),\quad
			\left(0,\left[\hat d_k/\hat a_k\right]_{[0,\gamma_{\max}]}\right),\\
			&\left(\beta_{\max},
			\left[(\hat d_k+\hat b_k\beta_{\max})/\hat a_k\right]_{[0,\gamma_{\max}]}\right),\\
			&\left(\left[(\hat e_k+\hat b_k\gamma_{\max})/\hat c_k\right]_{[0,\beta_{\max}]},
			\gamma_{\max}\right),
		\end{align*}
		where $[u]_{[l,r]}\triangleq\min\{r,\max\{l,u\}\}$.
		Consequently, its predicted decrease is at least as large as that of every single-direction or fixed-coupling comparator whose step-size pair lies in the same box.
	\end{proposition}
	
	The proposition provides both the analytical uncapped rule and the exact safeguarded box solution used by the implementation.
	
	\subsection{Model-Level Comparison and Algorithm Summary}\label{subsec:dominance}
	
	The closed-form rule makes the bi-direction update computationally explicit.
	Its model-level benefit is determined by comparing the optimized quadratic decrease with the decreases attainable after restricting the feasible step-size pair.
	
	\begin{proposition}[Predicted-decrease comparison]\label{prop:dominance}
		Consider first the uncapped controller in Proposition~\ref{prop:closed_form}, and fix a maximum comparator step size $s_{\max}>0$.
		Let $\widehat\Delta_k^{\bF}$ and $\widehat\Delta_k^{\bG}$ denote the best predicted decreases along the two coordinate axes, and let $\widehat\Delta_k^{\mathrm{2nd}}$ denote the best decrease along the PPM/EGM coupling $(\beta,\gamma)=(s,s^2)$:
		\[
		\begin{gathered}
			\widehat\Delta_{k}^{\bF}=\frac{[\hat e_k]_+^2}{2\hat c_k},
			\quad
			\widehat\Delta_{k}^{\bG}=\frac{[\hat d_k]_+^2}{2\hat a_k},\\
			\widehat\Delta_{k}^{\mathrm{2nd}}=-\min_{s\in[0,s_{\max}]}\hat q_k(s,s^2).
		\end{gathered}
		\]
		Then $\widehat\Delta_{k}^{\star}\ge\max\{\widehat\Delta_{k}^{\bF},\widehat\Delta_{k}^{\bG},\widehat\Delta_{k}^{\mathrm{2nd}}\}$ almost surely.
		Moreover, when the unconstrained critical point in Eq.~\eqref{eq:unconstrained_pair} belongs to $[0,\infty)^2$, the gap over the field-only rule is exactly
		\begin{equation}\label{eq:dominance_gap}
			\widehat\Delta_{k}^{\star}-\widehat\Delta_{k}^{\bF}
			=
			\frac{(\hat b_{k}\hat e_{k}+\hat c_{k}\hat d_{k})^{2}}{2\hat c_{k}\hat\delta_{k}}
			+\frac{[-\hat e_k]_{+}^{2}}{2\hat c_k}
			\ge 0,
		\end{equation}
		which is strictly positive if and only if either $\hat b_{k}\hat e_{k}+\hat c_{k}\hat d_{k}\neq 0$ or $\hat e_k<0$.
		In the usual field-descent regime $\hat e_k\ge0$, strict improvement is therefore equivalent to the curvature direction having a nonzero reduced gradient after projecting out the field direction in the $\widehat{\bH}_{k}$-metric.
	\end{proposition}
	
	Eq.~\eqref{eq:dominance_gap} measures the curvature direction's additional predicted decrease through the reduced gradient $\hat b_k\hat e_k+\hat c_k\hat d_k$ left after accounting for the field direction.
	For collinear directions, the common spectral shift makes the Hessian nonsingular; because the direction span is unchanged, any gain arises from step-size constraints or regularization.
	One joint policy update of the resulting controller consists of the following operations:
	\begin{enumerate}[(i)]
		\item obtain the stochastic field and curvature observations $(\hatF_k,\hatG_k)$;
		\item evaluate the five coefficients of the local Lyapunov model and apply the common spectral shift;
		\item solve the two-dimensional box QP for $(\beta_k,\gamma_k)$ using Proposition~\ref{prop:closed_form};
		\item apply the stochastic policy recursion in Eq.~\eqref{eq:update}.
	\end{enumerate}
	This sequence separates the geometric construction of the two search directions from their Lyapunov-based step-size selection.
	Section~\ref{sec:conv} combines this comparator dominance with the field-geometry certificate and explicitly propagates stochastic-oracle, model-estimation, Taylor-remainder, and safeguard errors into finite-time guarantees.
	
	% ============================================================
	\section{Finite-Time and Policy-Space Guarantees}\label{sec:conv}
	% ============================================================
	Section~\ref{sec:drift} constructs the controller and characterizes its predicted quadratic decrease.
	We now bound its realized and conditional Lyapunov changes, derive a finite-time certificate for the two available descent directions, and then identify when this certificate yields stationarity and policy-space performance guarantees.
	
	\subsection{Conditional Drift and a Bi-Direction Certificate}\label{subsec:pathwise}
	
	\begin{assumption}[Local admissible region]\label{ass:local}
		There exist an open set $\cN$ whose compact closure satisfies $\overline{\cN}\subset\cD$ and an index $k_{0}\ge0$ such that $\cN$ contains the local stationary component of interest, $\cP$ is $C^3$ on a neighborhood of $\overline{\cN}$, and $\bz_k\in\cN$ for every $k\ge k_0$ almost surely.
	\end{assumption}
	
	This is a local-invariance condition: the safeguards in Section~\ref{sec:impl} keep the accepted trajectory in a prescribed implementable region, and the analysis applies when that trajectory remains in $\cN$.
	Because $\overline{\cN}$ is compactly contained in $\cD$, continuity yields uniform constants $G_V,H_V,B_F,B_G\ge0$ such that $\norm{\nabla\Vk}\le G_V$, $\norm{\nabla^2\Vk}\le H_V$, $\norm{\bF}\le B_F$, and $\norm{\bG}\le B_G$ on $\cN$.
	In particular, $G_{V,k},H_{V,k}$ may be replaced by $G_V,H_V$ for $k\ge k_0$.
	
	Recall the pathwise reference quadratic $q_k$ in Eq.~\eqref{eq:qref} and the stabilized implemented quadratic $\hat q_k$ in Eq.~\eqref{eq:qhat}.
	For fixed caps $\beta_{\max},\gamma_{\max}>0$ and an implemented pair $(\beta_k,\gamma_k)\in[0,\beta_{\max}]\times[0,\gamma_{\max}]$, define the box-optimal predicted decrease
	\[
	\widehat\Delta_k^{\rm box}
	\triangleq
	-\min_{0\le\beta\le\beta_{\max},\,0\le\gamma\le\gamma_{\max}}
	\hat q_k(\beta,\gamma).
	\]
	The implemented model $\hat q_k$ uses the stabilized coefficients, and one-sided residuals measure model discrepancy and accepted step-size-pair suboptimality.
	Define the cubic Taylor remainder $\varepsilon_{k}^{\mathrm{rem}}\triangleq M_{3,k} (\beta_{k}\norm{\hatF_{k}}+\gamma_{k}\norm{\hatG_{k}})^{3}/6$.
	Let $\varepsilon_k^{\rm model}\ge0$ be $\mathscr F_{k+1}$-measurable and satisfy $q_k(\beta_k,\gamma_k)\le\hat q_k(\beta_k,\gamma_k) +\varepsilon_k^{\rm model}$.
	For the implemented box problem, define $\varepsilon_k^{\rm step,box}\triangleq\hat q_k(\beta_k,\gamma_k)+\widehat\Delta_k^{\rm box}\ge0$ and
	\begin{equation}\label{eq:eps_tot}
		\varepsilon_{k}^{\mathrm{box}}
		\triangleq
		\varepsilon_{k}^{\mathrm{rem}}+\varepsilon_{k}^{\mathrm{model}}
		+\varepsilon_{k}^{\mathrm{step,box}}.
	\end{equation}
	
	Suppose the raw evaluated coefficients satisfy $\lvert\tilde e_{k}-e_{k}\rvert\le\delta_{e,k}$, $\lvert\tilde d_{k}-d_{k}\rvert\le\delta_{d,k}$, $\lvert\tilde c_{k}-c_{k}\rvert\le\delta_{c,k}$, $\lvert\tilde a_{k}-a_{k}\rvert\le\delta_{a,k}$, and $\lvert\tilde b_{k}-b_{k}\rvert\le\delta_{b,k}$ for nonnegative error bounds.
	
	\begin{proposition}[Local-model and pathwise error bounds]\label{prop:model_approx}
		Under the preceding coefficient-error bounds, the following statements hold.
		\begin{enumerate}[(i)]
			\item For every feasible $(\beta,\gamma)$,
			\begin{equation}\label{eq:eps_model_bound}
				\begin{aligned}
					\varepsilon_k^{\rm model}(\beta,\gamma)
					&\triangleq\beta\delta_{e,k}+\gamma\delta_{d,k}
					+\tfrac12\beta^2\delta_{c,k}\\
					&\quad+\beta\gamma\delta_{b,k}+\tfrac12\gamma^2\delta_{a,k},\\
					\varepsilon_k^{\rm cmp}(\beta,\gamma)
					&\triangleq\varepsilon_k^{\rm model}(\beta,\gamma)
					+\tfrac12\rho_{c,k}\beta^2+\tfrac12\rho_{a,k}\gamma^2,\\
					q_k-\hat q_k&\le\varepsilon_k^{\rm model},\\
					\hat q_k-q_k&\le\varepsilon_k^{\rm cmp}.
				\end{aligned}
			\end{equation}
			The inequality arguments are suppressed, $\mathrm{cmp}$ denotes the comparator-side bound, and Eq.~\eqref{eq:eps_tot} uses $\varepsilon_k^{\mathrm{model}}=\varepsilon_k^{\mathrm{model}}(\beta_k,\gamma_k)$.
			\item Under Assumption~\ref{ass:local_reg}, the realized update satisfies
			\begin{equation}\label{eq:pathwise_drift}
				\Vk(\bz_{k+1})-\Vk(\bz_{k})
				\le -\widehat\Delta_{k}^{\rm box}+\varepsilon_{k}^{\rm box}
				\quad\text{almost surely}.
			\end{equation}
			\item If Assumptions~\ref{ass:oracles} and~\ref{ass:local_reg} hold and $\beta_{k}\le\beta_{\max}$, $\gamma_{k}\le\gamma_{\max}$, then
			\begin{equation}\label{eq:eps_rem_cond}
				\E[\varepsilon_{k}^{\mathrm{rem}}\mid\mathscr{F}_{k}]
				\le
				\frac{2M_{3,k}}{3}(\beta_{\max}+\gamma_{\max})^{3}\sigma_{3,k}^{3}.
			\end{equation}
			If $M_{3,k}\le\bar M_{3}$ and $\sigma_{3,k}\le\bar\sigma_{3}$, then $\E[\varepsilon_{k}^{\mathrm{rem}}\mid\mathscr F_k]\le\bar R_{\rm rem}$ almost surely, where $\bar R_{\rm rem}\triangleq(2\bar M_3/3)(\beta_{\max}+\gamma_{\max})^3\bar\sigma_3^3$.
		\end{enumerate}
	\end{proposition}
	
	The three parts control coefficient evaluation, pathwise model mismatch, and the conditional Taylor remainder; exact automatic differentiation and an exact box-QP solution set the corresponding evaluation, selected-pair, and solution residuals to zero.
	
	\begin{assumption}[Uniform local budgets]\label{ass:budget}
		There exist finite nonnegative constants appearing below such that, for every $k\ge k_0$, the following hold:
		(a) $\bar b_k^F\le\bar b_F$, $\bar b_k^G\le\bar b_G$, $\sigma_{2,k}\le\bar\sigma_2$, and $\sigma_{3,k}\le\bar\sigma_3$;
		(b) the common spectral shift satisfies $\rho_{c,k}=\rho_{a,k}\le\bar\rho$, and the raw coefficient-evaluation errors in Proposition~\ref{prop:model_approx} satisfy
		\[
		\begin{gathered}
			\E[\delta_{e,k}\mid\mathscr F_k]\le\bar\delta_e,
			\quad
			\E[\delta_{d,k}\mid\mathscr F_k]\le\bar\delta_d,\\
			\E[\delta_{c,k}\mid\mathscr F_k]\le\bar\delta_c,
			\quad
			\E[\delta_{a,k}\mid\mathscr F_k]\le\bar\delta_a,
			\quad
			\E[\delta_{b,k}\mid\mathscr F_k]\le\bar\delta_b;
		\end{gathered}
		\]
		(c) $M_{3,k}\le\bar M_3$ and $\E[\varepsilon_k^{\rm step,box}\mid\mathscr F_k]\le\bar R_{\rm step}$ almost surely.
	\end{assumption}
	
	Define the derived oracle-moment bounds and aggregate residual budgets
	\[
	\begin{gathered}
		M_F\triangleq3B_F^2+3\bar b_F^2+3\bar\sigma_2^2,
		\quad
		M_G\triangleq3B_G^2+3\bar b_G^2+3\bar\sigma_2^2,\\
		\bar b_D\triangleq\max\{\bar b_F,\bar b_G\},
		\quad
		M_D\triangleq\max\{M_F,M_G\},\\
		\bar\delta_1\triangleq\max\{\bar\delta_e,\bar\delta_d\},
		\quad
		\bar\delta_2\triangleq\max\{\bar\delta_c,\bar\delta_a\},\\
		\bar R_{\rm model}\triangleq{}
		\beta_{\max}\bar\delta_e+\gamma_{\max}\bar\delta_d
		+\tfrac12\beta_{\max}^2\bar\delta_c\\
		\quad{}+\beta_{\max}\gamma_{\max}\bar\delta_b
		+\tfrac12\gamma_{\max}^2\bar\delta_a,\\
		\bar R_{\rm box}\triangleq\bar R_{\rm rem}+\bar R_{\rm model}+\bar R_{\rm step}.
	\end{gathered}
	\]
	Proposition~\ref{prop:model_approx}, the step-size caps, and Assumption~\ref{ass:budget} give $\E[\varepsilon_k^{\rm box}\mid\mathscr F_k]\le\bar R_{\rm box}$ almost surely, so the aggregate is derived from its three displayed sources.
	For a comparator step size $s\in[0,\min\{\beta_{\max},\gamma_{\max}\}]$, let
	\begin{equation}\label{eq:R_direction}
		\mathcal R_D(s)
		\triangleq
		s(G_V\bar b_D+\bar\delta_1)
		+\tfrac12s^2(H_VM_D+\bar\rho+\bar\delta_2)
		+\bar R_{\rm box}.
	\end{equation}
	
	\begin{proposition}[Conditional and finite-time bi-direction Lyapunov certificate]\label{prop:directional_certificate}
		Suppose Assumptions~\ref{ass:smooth}, \ref{ass:oracles}, \ref{ass:local_reg}, \ref{ass:local}, and~\ref{ass:budget} hold.
		Then the following statements hold.
		\begin{enumerate}[(i)]
		\item For every $k\ge k_0$ and every admissible deterministic comparator step size $s$,
			\begin{equation}\label{eq:cond_drift}
				\mathscr D_k
				\le
				-s\left[\max\{D_F(\bz_k),D_G(\bz_k)\}\right]_+
				+\mathcal R_D(s).
			\end{equation}
			\item For every $s\in(0,\min\{\beta_{\max},\gamma_{\max}\}]$ and integer $K>k_0$,
			\begin{equation}\label{eq:directional_certificate}
				\begin{aligned}
					&\frac{1}{K-k_0}
					\sum_{k=k_0}^{K-1}
					\E\!\left[\left[\max\{D_F(\bz_k),D_G(\bz_k)\}\right]_+\right]
					\\
					&\quad\le
					\frac{\E[\Vk(\bz_{k_0})]}{s(K-k_0)}
					+\frac{\mathcal R_D(s)}{s}.
				\end{aligned}
			\end{equation}
		\end{enumerate}
	\end{proposition}
	
	Part~(i) transfers the strongest available single-direction decrease to the implemented update up to explicit residuals, and Part~(ii) converts this one-step comparison into a finite-time average certificate.
	The next subsection relates this signal to the stationarity residual.
	
	\subsection{From Geometric Coverage to Finite-Time Stationarity}\label{subsec:ergodic}
	
	For this conversion, define the geometric coverage modulus on $\cN$ by
	\begin{equation}\label{eq:coverage_modulus}
		m_{\rm geo}
		\triangleq
		\inf_{\substack{\bz\in\cN\\\Phi(\bz)>0}}
		\frac{\max\{\chi_F(\bz),\chi_G(\bz)\}}{\Phi(\bz)}.
	\end{equation}
	The condition $m_{\rm geo}>0$ states that, at every nonstationary point in $\cN$, at least one available direction supplies field-energy decrease proportional to the stationarity residual.
	The following proposition derives this property locally from the Jacobian at the target stationary point; $\sigma_{\min}(\cdot)$ denotes the smallest singular value.
	
	\begin{proposition}[A sufficient condition for geometric coverage]\label{prop:coverage_sufficient}
		Suppose Assumption~\ref{ass:smooth} holds.
		Fix a target stationary point $\bz^\star\in\mathcal Z^\star$ and define its metric-transformed Jacobian $\widetilde{\bA}_\star\triangleq\bW_V^{1/2}\bA(\bz^\star)\bW_V^{-1/2}$.
		If $\sym(\widetilde{\bA}_\star)$ is positive semidefinite and $\sigma_{\min}(\widetilde{\bA}_\star)>0$, then there exist a constant $c_\star>0$ and an open neighborhood $\cN_\star$ of $\bz^\star$, with $\overline{\cN_\star}\subset\cD$, such that
		\(\max\{\chi_F(\bz),\chi_G(\bz)\}\ge c_\star\Phi(\bz)\) for every $\bz\in\cN_\star$.
		Consequently, if the admissible region in Assumption~\ref{ass:local} satisfies $\overline{\cN}\subset\cN_\star$, then $m_{\rm geo}\ge c_\star>0$.
	\end{proposition}
	
	For $\bW_V=\Id$, max-variable concavity and min-variable convexity at $\bz^\star$ give the positive-semidefinite condition, while nonsingularity gives a nondegenerate target linearization; continuity then supplies neighborhood coverage, including pure rotation with $\sym(\widetilde{\bA}_\star)=0$.
	For the normal block $\bA=\mu\Id_2+\sigma J_0$, the exact ratio in Eq.~\eqref{eq:coverage_modulus} is $\max\{\mu,\sigma^2-\mu^2\}$; in particular, $\mu=0$ gives $m_{\rm geo}=\sigma^2$.
	
	Compactness and smoothness give the finite performance-direction bound
	\[
	B_{P,D}
	\triangleq
	\sup_{\bz\in\overline{\cN}}
	\max\!\left\{
	\left|\inner{\nabla\cP}{\bF}\right|,
	\left|\inner{\nabla\cP}{\bG}\right|
	\right\}.
	\]
	
	When Assumption~\ref{ass:local} holds and $m_{\rm geo}>0$, define the geometric descent coefficient $\mu_{\rm geo}$ and policy-component perturbation $r_P$ by
	\[
	\mu_{\rm geo}\triangleq\frac{\lambda_Fm_{\rm geo}}{\kappa_F},
	\quad
	r_P\triangleq\frac{\lambda_PB_{P,D}}{\kappa_P}.
	\]
	Selecting the larger signal in Eq.~\eqref{eq:composite_signals} and applying Eq.~\eqref{eq:coverage_modulus} gives, for every $\bz\in\cN$,
	\begin{equation}\label{eq:derived_drift_certificate}
		\max\{D_F(\bz),D_G(\bz)\}
		\ge
		\mu_{\rm geo}\Phi(\bz)-r_P.
	\end{equation}
	
	Proposition~\ref{prop:coverage_sufficient} and Eq.~\eqref{eq:derived_drift_certificate} provide the stationarity conversion: target-Jacobian conditions yield local coverage, which lower-bounds the strongest bi-direction Lyapunov signal by the stationarity residual up to $r_P$.
	The perturbation $r_P$ vanishes for $\lambda_P=0$ and otherwise enters the residual floor.
	
	\begin{theorem}[Main Theorem: Local Finite-Time Stationarity under Geometric Coverage]\label{thm:ergodic}
		Suppose Assumptions~\ref{ass:smooth}, \ref{ass:oracles}, \ref{ass:local_reg}, \ref{ass:local}, and~\ref{ass:budget} hold, and the modulus in Eq.~\eqref{eq:coverage_modulus} satisfies $m_{\rm geo}>0$.
		For any comparator step size $s\in(0,\min\{\beta_{\max},\gamma_{\max}\}]$, define $\mathcal R_{\rm geo}(s)\triangleq\mathcal R_D(s)+sr_P$.
		Then, for every integer $K>k_0$,
		\begin{equation}\label{eq:main_stationarity}
			\frac{1}{K-k_{0}}
			\sum_{k=k_{0}}^{K-1}\E[\Phi(\bz_{k})]
			\le
			\frac{\E[\Vk(\bz_{k_{0}})]}{s\mu_{\rm geo}(K-k_{0})}
			+
			\frac{\mathcal R_{\rm geo}(s)}{s\mu_{\rm geo}}.
		\end{equation}
	\end{theorem}
	
	For fixed admissible $s$ and $K$-independent budgets, Eq.~\eqref{eq:main_stationarity} gives an $O(1/K)$ transient to an explicit floor containing the oracle, model, Taylor-remainder, policy-component, and safeguard errors.
	The signals $\chi_F$ and $\chi_G$ respectively certify contractive symmetric fields and pure rotation, and box-QP dominance inherits either certificate.
	For an independent uniform index $I_K\in\{k_0,\ldots,K-1\}$, the left side is $\E[\Phi(\bz_{I_K})]$, which gives an equivalent uniform-random-iterate interpretation of the theorem.
	The comparator $s$ trades nominal decrease against higher-order and $s$-independent residuals, while the implemented controller continues to select $(\beta_k,\gamma_k)$ from the box QP.
	
	\subsection{Policy-Space Performance Guarantee}\label{subsec:performance_guarantee}
	
	Exploitability and robust return are policy-space quantities beyond parameter-space stationarity.
	We therefore transfer Lyapunov control to these performance measures.
	For stationary randomized policies, define the maximizing player's worst-case return, game value, and unregularized Nash gap (reported below as exploitability), where the subscript ${\rm BR}$ denotes best response:
	\begin{align}
		R_{\rm BR}(\pi)&\triangleq\inf_{\bar\nu}J(\pi,\bar\nu),\notag\\
		v^{\star}&\triangleq\sup_{\bar\pi}R_{\rm BR}(\bar\pi),\notag\\
		\operatorname{Gap}_{0}(\pi,\nu)
		&\triangleq
		\sup_{\bar\pi}J(\bar\pi,\nu)-\inf_{\bar\nu}J(\pi,\bar\nu).
		\label{eq:global_gap}
	\end{align}
	The following proposition links the entropy-regularized gap used by the controller to the original zero-sum Markov-game objective.
	
	\begin{proposition}[Policy-space performance bridge]\label{prop:performance_bridge}
		Suppose $\mathcal A$ and $\mathcal B$ are finite.
		For every policy pair $(\pi,\nu)$ and $\varepsilon\ge0$,
		\begin{equation}\label{eq:performance_bridge}
			\begin{aligned}
				0&\le v^{\star}-R_{\rm BR}(\pi)
				\le \operatorname{Gap}_{0}(\pi,\nu)\\
				&\le \operatorname{Gap}_{\varepsilon}(\pi,\nu)
				+\frac{\varepsilon\bigl(\log|\mathcal A|+\log|\mathcal B|\bigr)}{1-\alpha}.
			\end{aligned}
		\end{equation}
	\end{proposition}
	\begin{proof}
		The definitions in Eq.~\eqref{eq:global_gap} give the first inequality, while $v^{\star}\le\sup_{\bar\pi}J(\bar\pi,\nu)$ gives the second.
		The discounted entropy bounds are $0\le\mathsf H_\alpha^\pi\le\log|\mathcal A|/(1-\alpha)$ and $0\le\mathsf H_\alpha^\nu\le\log|\mathcal B|/(1-\alpha)$.
		Substituting them into the two best-response terms in Eq.~\eqref{eq:global_gap} gives the final inequality in Eq.~\eqref{eq:performance_bridge}.
	\end{proof}
	
	Thus the regularized gap controls the maximizing player's unregularized worst-case-return loss up to the stated entropy term.
	
	\begin{proposition}[Smooth finite-game performance component]\label{prop:smooth_entropy_gap}
		Suppose $\mathcal S,\mathcal A,\mathcal B$ are finite, the one-stage rewards are bounded, $\varepsilon>0$, and, on the parameter neighborhood under consideration, the current policies are $C^3$ with action probabilities bounded away from zero.
		Then the entropy-regularized best-response values against either current player, and hence $\operatorname{Gap}_\varepsilon$, are $C^3$ functions of the current policy parameters.
	\end{proposition}
	\begin{proof}
		For a fixed opponent policy, the finite-state soft Bellman operator is $C^3$ in its value vector and the opponent parameters because it is composed of smooth policy probabilities, finite sums, and log-sum-exp.
		Its value-vector Jacobian is a discounted transition matrix and therefore has induced infinity norm at most $\alpha<1$.
		The implicit-function theorem consequently makes its unique fixed point and regularized best-response value $C^3$ in the opponent parameters.
		Applying the same argument to the soft-min Bellman operator makes both best-response values, and hence their difference $\operatorname{Gap}_\varepsilon$, $C^3$.
	\end{proof}
	
	Consequently, choosing $\cP=\operatorname{Gap}_\varepsilon$ gives a concrete smooth performance component for Definition~\ref{def:Lyap} in finite Markov games.
	
	\begin{assumption}[Local policy-space error bound]\label{ass:performance_error_bound}
		There exist constants $c_{\rm gap}>0$ and $\xi_{\rm gap}\ge0$ such that
		\begin{equation}\label{eq:gap_error_bound}
			\operatorname{Gap}_\varepsilon(\bz)
			\le
			c_{\rm gap}\Phi(\bz)+\xi_{\rm gap},
			\quad
			\bz\in\cN.
		\end{equation}
	\end{assumption}
	
	Eq.~\eqref{eq:gap_error_bound} is a local Polyak--\L{}ojasiewicz-type policy error bound \cite{Karimi2016PL}.
	The slack $\xi_{\rm gap}$ accommodates parameter-stationary points with nonzero policy-space exploitability and hence allows the two zero sets to differ.
	
	\begin{corollary}[Ergodic policy-space performance]\label{cor:ergodic_performance}
		Suppose the assumptions of Theorem~\ref{thm:ergodic} and Assumption~\ref{ass:performance_error_bound} hold, and $\mathcal A$ and $\mathcal B$ are finite.
		Then
		\begin{equation}\label{eq:ergodic_performance}
			\begin{aligned}
				&\frac{1}{K-k_0}\sum_{k=k_0}^{K-1}
				\E[v^\star-R_{\rm BR}(\pi_{\btheta_k})]\\
				&\le
				\frac{1}{K-k_0}\sum_{k=k_0}^{K-1}
				\E[\operatorname{Gap}_0(\pi_{\btheta_k},\nu_{\bpsi_k})]\\
				&\le
				c_{\rm gap}\!\left[
				\frac{\E[\Vk(\bz_{k_0})]}{s\mu_{\rm geo}(K-k_0)}
				+\frac{\mathcal R_{\rm geo}(s)}{s\mu_{\rm geo}}
				\right]
				+\xi_{\rm gap}+B_\varepsilon,
			\end{aligned}
		\end{equation}
		where $B_\varepsilon\triangleq\varepsilon(\log|\mathcal A|+\log|\mathcal B|)/(1-\alpha)$ is the entropy-regularization transfer term.
	\end{corollary}
	\begin{proof}
		Average Eq.~\eqref{eq:gap_error_bound}, apply Theorem~\ref{thm:ergodic}, and then use Proposition~\ref{prop:performance_bridge} pointwise.
	\end{proof}
	
	Thus, Eq.~\eqref{eq:ergodic_performance} gives the policy-space consequence of the main stationarity theorem under the local error bound with slack.
	
	% ============================================================
	\section{Practical Implementation}\label{sec:impl}
	% ============================================================
	This section specifies how the local quantities required by the controller can be evaluated in practice.
	We consider exact automatic differentiation, directional-secant approximations, and safeguarded box-constrained updates.
	For finite games, the implementation evaluates both the saddle field and the regularized gap from $J_\varepsilon$ with $\varepsilon>0$; the unregularized return $J$ is retained for the performance measures in Section~\ref{subsec:performance_guarantee}.
	
	\subsection{Trajectory construction of the policy oracles}\label{subsec:trajectory_oracles}
	\textcolor{blue}{At update $k$, collect a batch $\mathcal B_k=\{\tau_{k,i}\}_{i=1}^{N_k}$ of $N_k$ trajectories with rollout length $L_k$ under the current policies and transition kernel; transition $t$ records $(s_{k,i,t},a_{k,i,t},b_{k,i,t},r_{k,i,t},s_{k,i,t+1})$.}
	\textcolor{blue}{For a candidate parameter $\bz=(\btheta,\bpsi)$, define the cumulative action log-likelihood and regularized stage payoff by}
	\[
	\textcolor{blue}{
	\begin{aligned}
	\ell_{k,i,t}(\bz)
	&\triangleq\sum_{u=0}^{t}\log\!\bigl[
	\pi_{\btheta}(a_{k,i,u}\mid s_{k,i,u})
	\nu_{\bpsi}(b_{k,i,u}\mid s_{k,i,u})\bigr],\\
	c_{\varepsilon,k,i,t}(\bz)
	&\triangleq r_{k,i,t}
	-\varepsilon\log\pi_{\btheta}(a_{k,i,t}\mid s_{k,i,t})\\[-2pt]
	&\quad
	+\varepsilon\log\nu_{\bpsi}(b_{k,i,t}\mid s_{k,i,t}).
	\end{aligned}}
	\]
\textcolor{blue}{Let $\operatorname{sg}$ preserve its argument in value but have zero derivative, and set $\mathcal M_{k,i,t}(\bz)\triangleq\exp\{\ell_{k,i,t}(\bz)-\operatorname{sg}(\ell_{k,i,t}(\bz))\}$.}
\textcolor{blue}{Using the policy-gradient score-function construction~\cite{Sutton2000PolicyGradient}, the infinitely differentiable Monte Carlo estimator (DiCE)~\cite{Foerster2018DiCE} gives the finite-trajectory surrogate}
	\[
	\textcolor{blue}{
	\widehat J_{\varepsilon,k}(\bz;\mathcal B_k)
	\triangleq
	\frac{1}{N_k}\sum_{i=1}^{N_k}\sum_{t=0}^{L_k-1}
	\alpha^t\mathcal M_{k,i,t}(\bz)c_{\varepsilon,k,i,t}(\bz).}
	\]
	\textcolor{blue}{Its first two automatic derivatives provide the sampled field and Jacobian--vector curvature estimate}
	\[
	\textcolor{blue}{
	\begin{aligned}
	\widehat{\bF}(\bz;\mathcal B_k)
	&\triangleq
	\begin{bmatrix}
	-\nabla_{\btheta}\widehat J_{\varepsilon,k}(\bz;\mathcal B_k)\\
	\phantom{-}\nabla_{\bpsi}\widehat J_{\varepsilon,k}(\bz;\mathcal B_k)
	\end{bmatrix},\\
	\hatF_k
	&\triangleq\widehat{\bF}(\bz_k;\mathcal B_k),\quad
	\hatG_k
	\triangleq
	\left.\nabla_{\bz}\widehat{\bF}(\bz;\mathcal B_k)\right|_{\bz=\bz_k}
	\operatorname{sg}(\hatF_k).
	\end{aligned}}
	\]
	\textcolor{blue}{Thus, automatic differentiation evaluates $\hatG_k$ without forming the Jacobian.}
	\begin{proposition}[Finite-trajectory oracle validity]\label{prop:trajectory_oracle_validity}
		\textcolor{blue}{Suppose Assumption~\ref{ass:smooth} holds; $N_k,L_k$ are finite positive $\mathscr F_k$-measurable integers; and rewards and the first two log-policy derivatives are uniformly bounded, with policy probabilities bounded below by $p_{\min}>0$, on the implementable region. Then $\hatF_k,\hatG_k$ are conditionally bounded, and their centered decomposition satisfies Assumption~\ref{ass:oracles} with finite $\mathscr F_k$-measurable bounds.}
	\end{proposition}
\textcolor{blue}{See \hyperref[app:trajectory_oracle_validity]{Appendix~\ref*{app:proofs}, item~i)} for the proof.}
	
	\subsection{Exact local-model evaluation}
	The five pathwise coefficients in Eq.~\eqref{eq:drift_terms_def} require $\hatF_k$, $\hatG_k$, $\nabla\Vk(\bz_k)$, and Hessian--vector products of $\Vk$ along the two directions. Automatic differentiation can pass through a fixed differentiable inner loop for $\cP$, while implicit differentiation handles an exact local envelope. Exact evaluation gives the reference coefficients, after which $\hat q_k-q_k=\lambda_k(\beta^2+\gamma^2)/2\ge0$.
	
	\subsection{Approximate local-model estimation}
	A directional-secant realization holds $(\hatF_k,\hatG_k)$ fixed and evaluates $\Vk$ at the axial points $\bz_k\pm\delta_{F,k}\hatF_k$, $\bz_k\pm\delta_{G,k}\hatG_k$ and the four mixed-sign points. Standard centered first-, second-, and mixed-difference stencils estimate $(e_k,d_k,c_k,a_k,b_k)$ before the spectral shift. With a common seed $\zeta_k$, the alternative curvature oracle $\hatG_k^{\rm fd}=[\hatF(\bz_k+\epsilon_k\hatF_k,\zeta_k)-\hatF(\bz_k,\zeta_k)]/\epsilon_k$ reduces sampling variation; exponential moving averages may smooth the quadratic coefficients. Proposition~\ref{prop:model_approx} quantifies the resulting model error.
	
	\subsection{Safeguarded computation}\label{subsec:safeguard}
	With positive caps $\beta_{\max}$ and $\gamma_{\max}$, the implementation solves $\min_{(\beta,\gamma)\in[0,\beta_{\max}]\times[0,\gamma_{\max}]}\hat q_{k}(\beta,\gamma)$.
	Any loss from approximate solution or backtracking relative to the box optimum is recorded in $\varepsilon_k^{\rm step,box}$ in Theorem~\ref{thm:ergodic}.
	Each iteration queries $(\hatF_k,\hatG_k)$, evaluates the five tilde coefficients, applies the spectral shift, solves the two-dimensional box QP by Proposition~\ref{prop:closed_form}, and applies Eq.~\eqref{eq:update}.
	If the trial segment leaves a prescribed implementable trust region $\mathcal T\subset\cD$, a common backtracking factor shrinks $(\beta_k,\gamma_k)$; the analysis assumes the accepted trajectory lies in $\cN$, and records lost model decrease in $\varepsilon_k^{\mathrm{step,box}}$.
	For bounded $\mathcal T$, the implementation reduces the probe radii as needed until every directional-stencil point lies in $\mathcal T$.
	
	% ============================================================
	\section{Experiments}\label{sec:experiments}
	% ============================================================
	
	The experiments first evaluate nonlinear neural-policy optimization with exact population directions and then replace those directions by sampled stochastic oracles while varying physical attack strength.
	This progression separates the effect of nonlinear policy parameterization from sampling error and environment severity.
	Both studies compare QP+G, its field-only ablation noG ($\gamma_k=0$), GDA, Adam-GDA \cite{KingmaBa2015Adam}, EGM, and the three-inner-iteration PPM approximation PPM-3.
	All methods update both players simultaneously; unless noted, fixed step sizes and QP caps are $0.03$, $\eta_\psi=1$ gives $\cE_F=\phi$, and population normalizers are frozen at initialization with floor $10^{-10}$.
	For every Markov-game experiment, $(\alpha,\varepsilon)=(0.9,0.03)$; the regularized and unregularized value tolerances are $10^{-10}$ and $10^{-12}$, and the spectral floor is $10^{-8}$.
	Population-oracle comparisons use paired two-sided Student-$t$ $95\%$ confidence intervals for metric-oriented differences, with positive values favoring QP+G; the finite-trajectory sweep reports means and standard errors.
	
\subsection{Neural-Policy Markov Games with Population Oracles}\label{subsec:neural_mg}
	The population-oracle experiment isolates nonlinear policy parameterization from sampling noise using separate $4$--$8$--$3$ tanh--softmax policies, forming a $134$-dimensional field on the four games specified in Appendix~\ref{app:environments}.
	The QP+G/noG controller uses the composite Lyapunov function
	\begin{equation}\label{eq:vic_merit}
		V_{\rm MG}(\bz)
		=0.3\frac{\tfrac12\norm{\bF(\bz)}^{2}}{\kappa_F}
		+\frac{\operatorname{Gap}_{\varepsilon}(\bz)}{\kappa_P},
	\end{equation}
	This corresponds to $\lambda_F=0.3$ and $\lambda_P=1$, with both normalizers frozen at initialization.
	An exact differentiable value-system solve supplies $\bF$ and $\bG=\nabla\bF\,\bF$; each stencil recomputes the regularized best responses, and QP+G and noG use nine- and three-point central stencils, respectively, with $\delta_D=\min\{10^{-2},10^{-3}/\max(\|D\|,10^{-12})\}$, $D\in\{\bF,\bG\}$.
	Five screening seeds select each baseline step size from $\{0.001,0.003,0.01,0.03\}$; ten disjoint test seeds use $60$ simultaneous updates.
	QP+G attains higher worst-case return and lower exploitability than noG in all four games, with positive lower endpoints in all paired $95\%$ confidence intervals for return gain.
	The largest exploitability reduction is $79.1\%$ on CyclicControl, and the nonzero curvature allocations across all four games show that the composite QP exploits gap-gradient and cross-curvature information beyond the field-energy criterion of Theorem~\ref{thm:skew_dominance}.
	Fig.~\ref{fig:neural_mg} reports the complete convergence trajectories for worst-case return, exploitability, regularized Nash gap, and field norm.

\begin{figure*}[!t]
\centering
\includegraphics[width=\textwidth]{fig_vi_c_neural.pdf}
\caption{Neural-policy population-oracle results over ten test seeds (mean $\pm$ standard error); boundary triangles mark Adam-GDA means outside the displayed range.}
\label{fig:neural_mg}
\end{figure*}
	
\subsection{\textcolor{blue}{Finite-Trajectory Robustness under Attack-Strength Shifts}}\label{subsec:finite_trajectory}
\textcolor{blue}{We next replace population directions by sampled oracles and vary the physical severity of the adverse interaction.}
\textcolor{blue}{In each neural-policy game, $\eta\in\{0.5,0.75,1,1.25,1.5,1.75,2\}$ scales only the adverse reward term; $\eta=1$ is nominal, while transitions, policies, initializations, and algorithmic parameters remain fixed.}
\textcolor{blue}{Specifically, $\eta$ scales negative cyclic outcomes, matched-action jamming loss, matched-route interdiction loss, and uncovered-target loss in CyclicControl, FrequencyHopping, RoutingInterdiction, and SecurityPatrol, respectively.}
\textcolor{blue}{The minimizing reward is its negative, preserving the zero-sum game.}

\textcolor{blue}{At each of $60$ updates, the two $4$--$8$--$3$ policies use one batch of $512$ horizon-$16$ trajectories ($8192$ transitions).}
\textcolor{blue}{The same detached-policy batch supplies the field, Hessian--vector curvature, and local-model stencils through unclipped per-decision prefix ratios and DiCE~\cite{Foerster2018DiCE}; population quantities are used only for evaluation.}
\textcolor{blue}{QP+G and noG use the independently normalized field-energy Lyapunov function with $\lambda_P=0$, making the sweep a direct test of stochastic field and curvature estimation.}
\textcolor{blue}{GDA, EGM, and PPM-3 use the common screened step size $0.03$, Adam-GDA uses $0.001$, and all step sizes are frozen before the five-seed severity sweep and are not retuned across $\eta$.}
\textcolor{blue}{EGM and PPM-3 reuse the same trajectory batch for their additional field queries, giving every method the same transition budget per update.}

\begin{figure*}[!t]
\centering
\includegraphics[width=\textwidth]{fig_vi_d_finite_trajectory_robustness_six_methods.pdf}
\caption{\textcolor{blue}{Finite-trajectory robustness over attack strength $\eta$ and five held-out seeds (mean $\pm$ standard error). The top row reports final unregularized worst-case return, for which higher is better; the bottom row reports exploitability, for which lower is better. The vertical dotted line marks the nominal setting $\eta=1$.}}
\label{fig:finite_trajectory_robustness}
\end{figure*}

\textcolor{blue}{Across all $28$ settings in Fig.~\ref{fig:finite_trajectory_robustness}, QP+G attains the lowest mean exploitability and the strongest overall mean worst-case-return profile, demonstrating curvature-enabled improvement under both finite-trajectory DiCE estimation and attack-strength shifts.}

	% ============================================================
	\section{Conclusion}
	% ============================================================
	We developed complementary field and curvature directions for zero-sum Markov-game policy optimization and a closed-form box QP that selects their step sizes from a second-order Lyapunov-drift model.
	The analysis identifies when curvature decreases field energy and, under local geometric coverage, gives an $O(1/K)$ ergodic transient plus an explicit stochastic-error floor, with transfers to Nash-gap and worst-case-return bounds.
	Neural population-oracle and finite-trajectory experiments show improved convergence and the strongest overall policy-space performance under attack-strength shifts.
	
\appendices
% ============================================================
	
	\section{Proofs}\label{app:proofs}
\paragraph{Proof of Proposition~\ref{prop:DTA}}\label{app:DTA}
	The GDA identity is immediate. On a closed ball $\bar B(\bz_k;r)\subset\cD$, let $M_i$ and $B_F$ bound $\nabla^i\bF$ ($i=1,2$) and $\bF$.
	Choose $\bar s_k>0$ such that $\bar s_k B_F\le r$, $\bar s_k M_1\le1/2$, and $\bar s_k\le1$.
	For $s\le\bar s_k$, $T_{k,s}(\bu)=\bz_k-s\bF(\bu)$ is a self-map and contraction, giving a unique PPM update with $\norm{\Delta_k^{\rm PPM}}\le sB_F$, where $\Delta_k^{\rm PPM}=\bz_{k+1}^{\rm PPM}-\bz_k$; the EGM points likewise remain in the ball.
	Taylor expansion of the PPM equation gives $\Delta_k^{\rm PPM}+s\bF(\bz_k)=-s\nabla\bF(\bz_k)\Delta_k^{\rm PPM}-s\mathbf r_{2,k}$, where $\norm{\mathbf r_{2,k}}\le M_2\norm{\Delta_k^{\rm PPM}}^2/2$.
	Thus $\norm{\Delta_k^{\rm PPM}+s\bF(\bz_k)}\le s^2M_1B_F+s^3M_2B_F^2/2$; substituting once more yields the stated PPM expansion and $\norm{R_s^{\rm PPM}(\bz_k)}\le C_{\rm PPM}s^3$.
	For EGM, Taylor's theorem gives $\bF(\bz_k-s\bF(\bz_k))=\bF(\bz_k)-s\nabla\bF(\bz_k)\bF(\bz_k)+\mathbf r_{E,k}$ with $\norm{\mathbf r_{E,k}}\le M_2s^2B_F^2/2$; multiplication by $-s$ yields the stated expansion with $\norm{R_s^{\rm EGM}(\bz_k)}\le(M_2B_F^2/2)s^3$.
	\hfill$\blacksquare$
	
	\paragraph{Proof of Theorem~\ref{thm:skew_dominance}}\label{app:skew_dominance}
	With $\bA=\nabla\bF$ and $\nabla\phi=\bA^\top\bF$, direct substitution gives Eq.~\eqref{eq:dphi_F} and $\inner{\nabla\phi}{\bG}=\bF^\top\bA^2\bF$.
	Expanding $\bA^2$ as $\bS^2+\bS\bW+\bW\bS+\bW^2$ eliminates the antisymmetric cross term, while $\bF^\top\bS^2\bF=\norm{\bS\bF}^2$ and $\bF^\top\bW^2\bF=-\norm{\bW\bF}^2$, proving Eq.~\eqref{eq:dphi_G}.
	The right-hand side of Eq.~\eqref{eq:dphi_G} is negative exactly under Eq.~\eqref{eq:skew_dominance}.
	\hfill$\blacksquare$
	
	\paragraph{Proof of Proposition~\ref{prop:closed_form}}\label{app:closed_form}
	Positive definiteness makes $\hat q_k$ strictly convex. Solving $\widehat{\bH}_k(\beta,\gamma)^\top=(\hat e_k,\hat d_k)^\top$ and inverting gives Eqs.~\eqref{eq:unconstrained_pair}--\eqref{eq:Delta_hat_interior}. Uncapped boundary solutions lie on the axes; capped ones lie on four edges, where scalar minimization and clipping give the displayed candidates. Strict convexity and objective comparison prove uniqueness and the comparator claim.
	\hfill$\blacksquare$
	
	\paragraph{Proof of Proposition~\ref{prop:dominance}}\label{app:dominance}
	The comparator sets are subsets of $[0,\infty)^2$, so set inclusion and scalar minimization prove dominance and their formulas.
	Completing the square in Eq.~\eqref{eq:Delta_hat_interior} gives $\widehat\Delta_k^\star=\hat e_k^2/(2\hat c_k)+ (\hat b_k\hat e_k+\hat c_k\hat d_k)^2/(2\hat c_k\hat\delta_k)$.
	Subtracting $[\hat e_k]_+^2/(2\hat c_k)$ proves Eq.~\eqref{eq:dominance_gap} and its equality condition.
	\hfill$\blacksquare$
	
	\paragraph{Proof of Proposition~\ref{prop:coverage_sufficient}}\label{app:coverage_sufficient}
	Let $\widetilde{\bA}=\bW_V^{1/2}\bA\bW_V^{-1/2}$, $\widetilde{\bS}=\sym(\widetilde{\bA})$, $\widetilde{\bW}=\skewop(\widetilde{\bA})$, and $h(\bz,\bu)=\max\{\bu^\top\widetilde{\bS}\bu,\norm{\widetilde{\bW}\bu}^2-\norm{\widetilde{\bS}\bu}^2\}$.
	At $\bz^\star$, every unit $\bu$ has either $\bu^\top\widetilde{\bS}\bu>0$, or, by positive semidefiniteness, $\widetilde{\bS}\bu=0$ and hence $\norm{\widetilde{\bW}\bu}=\norm{\widetilde{\bA}\bu}>0$ by nonsingularity. Compactness and continuity give $h\ge c_\star>0$ uniformly on the unit sphere in a neighborhood of $\bz^\star$. For $\Phi>0$, substituting $\bu=\bW_V^{1/2}\bF/\sqrt{\Phi}$ gives $\chi_F/\Phi=\bu^\top\widetilde{\bS}\bu$ and $\chi_G/\Phi=\norm{\widetilde{\bW}\bu}^{2}-\norm{\widetilde{\bS}\bu}^{2}$.
	\hfill$\blacksquare$
	
	\paragraph{Proof of Proposition~\ref{prop:model_approx}}\label{app:model_approx}
	Subtracting the two models, using $\hat c_k=\tilde c_k+\rho_{c,k}$ and $\hat a_k=\tilde a_k+\rho_{a,k}$, and dropping the nonpositive shifts in $q_k-\hat q_k$ gives the first bound in Eq.~\eqref{eq:eps_model_bound}; reversing the difference gives the comparator bound. Eq.~\eqref{eq:qref}, Taylor's bound, and the definitions of $\varepsilon_k^{\rm model}$ and $\varepsilon_k^{\rm step,box}$ then give Eq.~\eqref{eq:pathwise_drift}. Finally, the caps and $(a+b)^3\le4(a^3+b^3)$ yield
	$\varepsilon_k^{\rm rem}\le(2M_{3,k}/3)(\beta_{\max}+\gamma_{\max})^3(\norm{\hatF_k}^3+\norm{\hatG_k}^3)$; conditional expectation proves Eq.~\eqref{eq:eps_rem_cond} and its uniform specialization.
	\hfill$\blacksquare$
	
	\paragraph{Proof of Proposition~\ref{prop:directional_certificate}}\label{app:directional_certificate}
	For $k\ge k_0$, compare the box-QP solution with $(s,0)$ or $(0,s)$ according to the larger positive signal, and with $(0,0)$ otherwise; deterministic ties make this $\mathscr F_k$-measurable. Box optimality and Proposition~\ref{prop:model_approx} bound the realized increment by this comparator plus the coefficient and box-solution errors. Its conditional linear term is at most $-s[\max\{D_F,D_G\}]_++sG_V\bar b_D$, the moment bounds control the Hessian terms, and the zero comparator gives $\mathscr D_k\le\bar R_{\rm box}$. Collecting terms gives Eq.~\eqref{eq:cond_drift}; summation and $\Vk\ge0$ give Eq.~\eqref{eq:directional_certificate}.
	\hfill$\blacksquare$
	
	\paragraph{Proof of Theorem~\ref{thm:ergodic}}\label{app:ergodic}
	Let $S_k\triangleq[\max\{D_F(\bz_k),D_G(\bz_k)\}]_+$.
	Eq.~\eqref{eq:derived_drift_certificate} gives $S_k\ge\mu_{\rm geo}\Phi(\bz_k)-r_P$.
	Taking expectations and averaging, then applying Proposition~\ref{prop:directional_certificate}, yields $\mu_{\rm geo}(K-k_0)^{-1}\sum_{k=k_0}^{K-1}\E[\Phi(\bz_k)]\le \E[\Vk(\bz_{k_0})]/[s(K-k_0)]+\mathcal R_D(s)/s+r_P$.
	Division by $\mu_{\rm geo}>0$ and $\mathcal R_{\rm geo}(s)=\mathcal R_D(s)+sr_P$ give Eq.~\eqref{eq:main_stationarity}.
	\hfill$\blacksquare$

	\paragraph{\textcolor{blue}{Proof of Proposition~\ref{prop:trajectory_oracle_validity}}}\label{app:trajectory_oracle_validity}
	\textcolor{blue}{Stop-gradient differentiation gives $\mathcal M=1$, $\nabla\mathcal M=\nabla\ell$, and $\nabla^2\mathcal M=\nabla^2\ell+\nabla\ell\nabla\ell^{\top}$. Finite $L_k$ and the stated bounds therefore bound $\hatF_k,\hatG_k$ and their population counterparts. The centered decompositions have zero-mean noises, finite bias norms, and finite conditional second and third moments, which supply the bounds in Assumption~\ref{ass:oracles}.}
	\hfill\textcolor{blue}{$\blacksquare$}
	
\section{Evaluated Finite Markov Games}\label{app:environments}
All four games use actions $a,b\in\{0,1,2\}$ and a uniform initial state.
For $S>1$, the transition law is $P(s'\mid s,a,b)=(1-\xi)\mathbf 1_{\{s'=T(s,a,b)\}}+\xi\mathbf 1_{\{s'\ne T(s,a,b)\}}/(S-1)$, where $T$ is the nominal next-state map and $\xi$ is the transition-noise probability.
Let $\varphi_s=2\pi s/S$.
Let $M=\left[\begin{smallmatrix}0&-1&1\\1&0&-1\\-1&1&0\end{smallmatrix}\right]$.
In \emph{CyclicControl}, the players select opposing control modes; $S=5$, $\xi=0.02$, $T=(s+1+((a-b)\bmod3))\bmod S$, and $r_s(a,b)=(1+0.12\cos\varphi_s)M_{\pi_s(a),\pi_s(b)}$, where $\pi_s(j)=(j-s)\bmod3$.
In \emph{FrequencyHopping}, a transmitter selects frequency $a$ while a jammer selects frequency $b$; $S=8$, $\xi=0.08$, $T=(s+1+a-b)\bmod S$, and $q_s(a)=0.25\cos(\varphi_s+2\pi a/3)$.
Its reward is $q_s(a)+0.55$ for $a\ne b$ and $q_s(a)-1.10$ for $a=b$.
In \emph{RoutingInterdiction}, a routing player selects route $a$ while an interdictor selects route $b$; $S=6$, $\xi=0.10$, $T=(s+1+\mathbf 1_{\{a=b\}}+a)\bmod S$, and $c_s(a)=1+0.20\sin(\varphi_s+2\pi a/3)$.
Its reward is $c_s(a)-0.55-1.35\mathbf 1_{\{a=b\}}-0.15\mathbf 1_{\{a\ne b\}}$.
In \emph{SecurityPatrol}, a defender selects patrol target $a$ while an attacker selects target $b$; $S=8$, $\xi=0.08$, $T=(s+1+b+\mathbf 1_{\{a\ne b\}})\bmod S$, and $v_s(b)=0.8+0.25\cos(\varphi_s+2\pi b/3)$.
Its reward is $v_s(b)$ when $a=b$ and $-0.55v_s(b)$ otherwise.
The robustness sweep scales the negative cyclic outcomes, matched-channel jamming loss, matched-route interdiction loss, and uncovered-target loss by $\eta$ in the four games, respectively.
At every $\eta$, the minimizing player's reward is the negative of the maximizing player's reward, preserving the zero-sum formulation.
Neural policies use $x_s=(\sin\varphi_s,\cos\varphi_s,u_s,1)$, where $u_s$ is a normalized state or environment statistic.

\begingroup
\emergencystretch=2em
	\bibliographystyle{IEEEtran}
	\bibliography{refs}
	\endgroup
\end{document}
